\documentclass[JEP,XML,SOM,Unicode,NoFloatCountersInSection,published]{cedram}
\datereceived{2021-03-23}
\dateaccepted{2022-06-28}
\dateepreuves{2022-07-06}

\usepackage{mathrsfs}
\let\mathcal\mathscr
\multlinegap0pt
\def\mto{\mathchoice{\longmapsto}{\mapsto}{\mapsto}{\mapsto}}
\newcommand{\Psfrac}[2]{(\sfrac{#1}{#2})}
\newcommand{\psfrac}[2]{\sfrac{(#1)}{#2}}
\newcommand{\spfrac}[2]{\sfrac{#1}{(#2)}}
\newcommand{\pspfrac}[2]{\sfrac{(#1)}{(#2)}}
\newcommand{\bnorm}[1]{\bigl\| #1 \bigr\|}
\newcommand{\nnorm}[1]{\| #1\|}
\DeclareMathOperator{\osc}{osc}

\newcommand{\Changel}{\mathcode`l="7160}
\newcommand{\Changelback}{\mathcode`l="716C}
\newcommand{\NoChangel}[1]{%
\expandafter\let\csname old\string#1\endcsname=#1
\let#1=\relax
\newcommand{#1}{\mathcode`l="716C\csname old\string#1\endcsname\mathcode`l="7160 }%
}

\NoChangel{\log}
\NoChangel{\ln}
\NoChangel{\lim}
\NoChangel{\limsup}
\NoChangel{\liminf}
\NoChangel{\varlimsup}
\NoChangel{\varliminf}
\NoChangel{\varinjlim}
\NoChangel{\varprojlim}


\makeatletter
\def\@settitle{%
\vspace*{-8mm}
\raggedleft\includegraphics[scale=.5]{titre-jep}
\vtop to 50 mm{%
 \parindent=0pt
 {\abstractfont\article@logo\par}
 \medskip
 \hrule
 \vfil
 \begin{center}
 \def\baselinestretch{1.2}\large\vfil
   {\didottitraille\MakeUppercase\@title\par}
 \vfil\vfil
 \begin{minipage}{.8\textwidth}\centering
   \ifx\@empty\smfbyname\else
   {\smf@byfont\smfbyname\ifsmf@byauthor\enspace\else\ \fi}%
   \fi {\smf@authorfont \edef\smfandname{{\noexpand\smf@andfont
         \smfandname}} \andify\authors\authors\par}
 \end{minipage}
 \vfil \vrule height .4pt width .3\textwidth \vfil
 \end{center}}%
 \par\enlargethispage{.5\baselineskip}%
}
\def\@setthanks{\def\thanks##1{\par##1\@addpunct{{\upshape.}}}\vspace*{-5pt}\thankses}
\makeatother

\usepackage{mathtools}
\usepackage{esint}

\newtheorem{theorem}{Theorem}
\newtheorem{proposition}[theorem]{Proposition}
\newtheorem{lemma}[theorem]{Lemma}
\newtheorem{corollary}[theorem]{Corollary}
\theoremstyle{definition}
\newtheorem{definition}[theorem]{Definition}
\newtheorem{conjecture}[theorem]{Conjecture}
\theoremstyle{remark}
\newtheorem{remark}[theorem]{Remark}
\newtheorem{remarks}[theorem]{Remarks}
\newtheorem{example}[theorem]{Example}

\newcommand{\norm}[1]{\left\| #1 \right\|}
\newcommand\N{{\mathbb N}}
\newcommand\R{{\mathbb R}}
\newcommand\T{{\mathbb T}}
\newcommand\CC{{\mathbb C}}\let\C\CC
\newcommand\Q{{\mathbb Q}}
\newcommand\Z{{\mathbb Z}}
\newcommand{\cC}{\mathcal C}
\newcommand{\cQ}{\mathcal Q}
\newcommand{\cU}{\mathcal{U}}
\newcommand{\cK}{\mathcal K}
\newcommand{\cT}{\mathcal T}
\newcommand{\dd}{{\, \mathrm d}}
\let\var\varepsilon

\datepublished{2022-07-13}
\begin{document}
\frontmatter

\title[Quantitative De~Giorgi methods in kinetic
theory]{Quantitative De~Giorgi methods\\ in kinetic theory}

\author[\initial{J.} \lastname{Guerand}]{\firstname{Jessica} \lastname{Guerand}}
\address{Université de Montpellier, IMAG\\
499-554 rue du Truel, 34090 Montpellier, France}
\email{jessica.guerand@umontpellier.fr}
\urladdr{https://sites.google.com/view/guerand}

\author[\initial{C.} \lastname{Mouhot}]{\firstname{Clément} \lastname{Mouhot}}
\address{University of Cambridge, Department of
Pure Mathematics and Mathematical Statistics\\
Wilberforce Road,
Cambridge CB3 0WA, United Kingdom}
\email{c.mouhot@dpmms.cam.ac.uk}
\urladdr{https://cmouhot.wordpress.com/}

\thanks{The authors acknowledge funding by the ERC grant MAFRAN 2017-2022}

\begin{abstract}
We consider hypoelliptic equations of kinetic Fokker-Planck type, also known as Kolmogorov or ultraparabolic equations, with rough coefficients in the drift-diffusion operator. We give novel short quantitative proofs of the De~Giorgi intermediate-value Lemma as well as weak Harnack and Harnack inequalities. This implies Hölder continuity with quantitative estimates. The paper is self-contained.
\end{abstract}

\subjclass{35K70, 35Q84, 35R09, 35B45, 35B65}

\keywords{Hypoelliptic equations, kinetic theory, Fokker-Planck equation, ultraparabolic equations, Kolmogorov equation, Hölder continuity, De~Giorgi method, Moser iteration, averaging lemma, weak Harnack inequality, trajectories}

\altkeywords{Équations hypoelliptiques, théorie cinétique, équation de Fokker-Planck, équations ultraparaboliques, équation de Kolmogorov, continuité höldérienne, méthode de De~Giorgi, itération de Moser, lemme de moyenne, inégalité de Harnack faible, trajectoires}

\alttitle{Méthodes à la De~Giorgi quantitatives en théorie cinétique}

\begin{altabstract}
Nous considérons des équations hypoelliptiques de type Fokker-Planck cinétique, également appelées équations de Kolmogorov ou ultraparaboliques, avec des coefficients sans régularité dans l'opérateur de dérive-diffusion. Nous donnons de nouvelles preuves quantitatives du lemme des valeurs inter\-médiaires de De~Giorgi ainsi que des inégalités de Harnack faibles et fortes. Cela implique la continuité höldérienne avec bornes explicites. L'article ne fait pas appel à des résultats précédents.\end{altabstract}

\maketitle
\vspace*{-\baselineskip}
\tableofcontents
\mainmatter

\section{Introduction}

\label{sec:intro}

\subsection{The problem studied}
\label{sec:theequation}

This paper is concerned with local regularity properties, namely
boundedness, Harnack inequalities and Hölder continuity, of
solutions $f=f(t,x,v)$ to the following class of hypoelliptic
partial differential equations in divergence form
\begin{equation}
\label{e:main}
\partial_t f + v \cdot \nabla_x f = \nabla_v \cdot (
A \nabla_v f) + B \cdot \nabla_v f + S,
\quad t \in \R, \ x \in \R^d, \ v \in \R^d,
\end{equation}
where $A=A(t,x,v)$, $B=B(t,x,v)$ and $S=S(t,x,v)$ satisfy (for some
constants $0 < \lambda < \Lambda$):
\begin{equation}
\label{e:hyp-coef}
\begin{dcases}
\text{$A$ is a measurable symmetric real matrix field with
eigenvalues in } [\lambda,\Lambda],\\
\text{$B$ is a measurable vector field such that }
|B| \le \Lambda,\\
\text{$S$ is a real scalar field in $L^\infty$.}
\end{dcases}
\end{equation}
This equation naturally appears in kinetic theory where it is
refereed to as the \emph{kinetic Fokker-Planck equation}; it is
included in the class considered by Kolmogorov~\cite{kolmogorov}
(with constant $A$ and linear $B$) that inspired the theory of
hypoellipticity of Hörmander~\cite{MR0222474} (see
\cite{AP20}). The coefficients are called ``rough'' because $A$,
$B$ and $S$ in the drift-diffusion operator on the $v$ variable
are merely measurable with no further regularity.

Our class~\eqref{e:main}--\eqref{e:hyp-coef} is invariant under
translations in $t$, $x$ and under \emph{Galilean translations},
\ie under $z \mto z_0 \circ z$ where $z_0 =(t_0,x_0,v_0)$,
$z = (t,x,v)$ and with the non-commutative group operation
\begin{equation*}
z_0 \circ z = (t_0+t,x_0+x + t v_0, v_0 + v).
\end{equation*}
Finally for any $r>0$ it is invariant under the scaling
$z=(t,x,v) \mto rz := (r^2t,r^3x,r v)$. Using the invariances of
the equation, we define for $z_0 \in \R^{1+2d}$ and $r>0$:
\begin{align*}
Q_r(z_0) := z_0 \circ \left[ r Q_1 \right]
:= \left\{ -r^2 <t-t_0 \le 0, \ |x - x_0- (t-t_0)v_0|
< r^3, \ |v-v_0| < r \right\}
\end{align*}
and we simply write $Q_r(0)=Q_r$ when $z_0=0$. We denote $|E|$
the Lebesgue measure of a Lebesgue set $E$. We write
$a \lesssim b$ (resp. $a \gtrsim b$) when $a \le Cb$ (resp.
$a \ge Cb$) for some constant $C>0$ whose only relevant
dependency, if any, is specified in the index, as in
$\lesssim_{\text{parameter}}$. We write $a\sim b$ if
$a \lesssim b$ and $a \gtrsim b$. We write $\fint$ for integrals
normalized by the volume of the integration domain, and
$\cT:=\partial_t+v\cdot \nabla_x$.

\skpt
\begin{definition}[Weak solution, sub-solution, super-solution]
\label{d:weak}
Let $\cU = (a,b) \times \Omega_x \times \Omega_v$ with
$-\infty<a < b \le +\infty$ and $\Omega_x$ and $\Omega_v$ two
open sets of $\R^d$. A function $f: \cU \to \R$ is a \emph{weak
solution} of \eqref{e:main} on $\cU$ if
\[
f \in L^\infty ((a,b);L^2 (\Omega_x \times \Omega_v)) \cap L^2
((a,b) \times \Omega_x;H^1(\Omega_v))
\]
and \eqref{e:main} is
satisfied in the sense of distributions in $\cU$. A function
$f$ is a \emph{weak sub-solution} of \eqref{e:main} if
\[
f \in L^\infty ((a,b);L^2 (\Omega_x \times \Omega_v)) \cap L^2
((a,b) \times \Omega_x;H^1(\Omega_v))
\]
and for all
$\beta : \R \to \R$ in $C^2$ with $\beta' \ge 0$ and
$\beta'' \ge 0$ both bounded, and any non-negative
$\varphi \in C^\infty_c (\cU)$,
\begin{equation*}
- \int_\cU \beta(f) \cT \varphi \dd z
\le - \int_{\cU} A \nabla_v \beta(f) \cdot \nabla_v \varphi \dd
z + \int_{\cU} \left[ B \cdot \nabla_v \beta(f)+S \beta'(f)
\right] \varphi \dd z.
\end{equation*}
It is a \emph{weak super-solution} of \eqref{e:main} if $-f$ is
a weak sub-solution.
\end{definition}

\begin{remark}
\label{r:def gen}
This definition is equivalent to those in \cite{pp} and
\cite{gimv} in the case of solutions, but is weaker than them
in the case of sub- and super-solutions. Indeed~\cite{pp,gimv}
make respectively the extra regularity assumption
$\cT f \in L^{2}((a,b) \times \Omega_x \times \Omega_v)$ or
$\cT f \in L^2 ((a,b) \times \Omega_x;H^{-1}(\Omega_v))$.
These assumptions were introduced to justify the energy
estimates. It is however enough to assume the renormalization
formulation above, and it allows to include important
sub-solutions such as for instance $f=f(t)={\bf 1}_{t \le 0}$
(when $S=0$) which were excluded by the definition
in~\cite{pp,gimv}. Our definition is equivalent to that of De
Giorgi in the elliptic case (and reminiscent of the definition
of solutions in \cite{GV15}).
\end{remark}

\subsection{Main contributions}
\label{sec:contributions}

Given the invariances, we only state results in unit
centered cylinders.
\begin{figure}[htb]
\includegraphics[scale =.45]{LVI.pdf}
\caption{The different cylinders in the Intermediate-Value
Lemma and Harnack inequalities.}
\label{fig:LVI}
\end{figure}
\begin{theorem}[Intermediate-Value Lemma]
\label{t:IVL}
Given $\delta_1,\delta_2 \in (0,1)$, there are explicit
constants
$r_0 \sim
\sfrac{\sqrt{\delta_1}}{\sqrt{1+\|S\|_{L^\infty(Q_1)}}}$ in
$(0,\sfrac{1}{20})$ if $S \not =0$ and $r_0=\sfrac{1}{20}$ if
$S=0$, and
$\theta \sim
\sfrac{(\delta_1\delta_2)^{10d+15}}{(1+\|S\|_{L^\infty(Q_1)})^{4d+3}}$
and
$\nu \gtrsim \sfrac{\left( \delta_1 \delta_2
\right)^{10d+16}}{(1+\|S\|_{L^\infty(Q_1)})^{2d+1}}$ both in
$(0,1)$, such that any sub-solution
$f:Q_1\to \mathbb{R}$ to
\eqref{e:main}--\eqref{e:hyp-coef} so that $f\leq 1$ in
$Q_{\sfrac{1}{2}}$ and
\begin{equation}
\label{e: hyp IVL}
|\{f\leq 0\}\cap Q_{r_0}^{-}| \geq \delta_1
|Q_{r_0}^{-}|
\quad \text{ and } \quad
|\{f\geq 1-\theta\} \cap Q_{r_0}| \geq \delta_2
|Q_{r_0}|,
\end{equation}
\ie we control the measure of where $f$ is below $0$ and above
$(1-\theta)$, satisfies
\begin{equation}
\label{e: concl IVL}
\left|\left\{ 0<f< 1-\theta \right\}\cap Q_{\sfrac12} \right|
\geq \nu |Q_{\sfrac12}|,
\end{equation}
where
$Q_{r_0}^{-} := Q_{r_0}(-2r_0^2,0,0) =(-3r_0^2,-2r_0^2]
\times B_{r_0^3} \times B_{r_0}$ (see Figure~\ref{fig:LVI}).
\end{theorem}

\begin{remark}
\label{rem:gap}
This lemma is the kinetic quantitative counterpart of the
quantitative elliptic \cite{DG56,DG57,vasseur} and parabolic
\cite{gueDGhalv2} intermediate value lemma. As~in the
parabolic case, past and a future cylinders $Q_{r_0}^{-}$ and
$Q_{r_0}^{+}$ are required to be disjoint but contrary to the
parabolic case, a gap in time between the two cylinders is also
required. This gap is also mentioned in~\cite{gimv,AP20}. Let
us explain why it cannot be removed. Consider for instance
$S=0$ and velocities bounded by $|v| \le V_m$ in the
cylinder. Then ${\bf 1}_{x+ct<a}$ is a sub-solution for any
$a\in \mathbb{R}$ and $|c|>V_m$. If $Q_{r_0}^{-}$ and $Q_{r_0}$
were too close, a line of discontinuity of the form $x+ct=a$
could cross both and the previous sub-solution would contradict
the conclusion of Theorem~\ref{t:IVL}.
\end{remark}

\begin{theorem}[Harnack inequalities]
\label{t:harnack}
There is $\zeta >0$ depending only $\lambda,\Lambda$ such that
any non-negative weak super-solution $f$
to~\eqref{e:main}--\eqref{e:hyp-coef} in $Q_1$ satisfies the
weak Harnack inequality\vspace*{-3pt}
\begin{equation}
\label{eq:w-Harnack-stat}
\biggl( \int_{\tilde Q_{\sfrac{r_0}{2}}^-} f^\zeta (z) \dd t \dd x \dd v
\biggr)^{\sfrac{1}{\zeta}}
\lesssim_{\lambda,\Lambda} \inf_{ Q_{\sfrac{r_0}{2}}}f
+ \|S\|_{L^\infty (Q_1)},
\end{equation}
where $r_0=\sfrac{1}{20}$ and
$\tilde Q_{\sfrac{r_0}{2}}
^{-}:=Q_{\sfrac{r_0}{2}}((-\frac{19}{8}\,r_0^2,0,0))$ (see
Figure~\ref{fig:LVI}), and any non-negative weak solution $f$
to~\eqref{e:main}--\eqref{e:hyp-coef} satisfies the following
Harnack inequality (with
$\tilde
Q_{\sfrac{r_0}{4}}^{-}:=Q_{\sfrac{r_0}{4}}((-\frac{19}{8}\,r_0^2,0,0))$)\vspace*{-3pt}
\begin{equation}
\label{eq:s-Harnack-stat}
\sup_{\tilde Q_{\sfrac{r_0}{4}}^{-}} f \lesssim_{\lambda,\Lambda}
\inf_{Q_{\sfrac{r_0}{4}}} f + \| S \|_{L^\infty(Q_1)}.
\end{equation}
\end{theorem}

\skpt
\begin{remarks}
\begin{enumerate}
\item The ``weak'' Harnack inequality, in spite of its name, is
not weaker than Harnack inequality since it holds for
super-solutions. Combined with the $L^{\zeta} \to L^\infty$
gain of integrability in~Proposition~\ref{prop:1st lemma}, it
implies the Harnack inequality for solutions. Super-solutions
of the form ${\bf 1}_{x+ct\ge a}$ for $a\in \mathbb{R}$ and
$|c|>V_m$ (included in our definition) show that the gap in
time is required in~\eqref{eq:w-Harnack-stat}.
\item The Harnack inequality for equation~\eqref{e:main} was
first proved in~\cite{gimv} by a non-constructive argument. The
present paper provides a new constructive De~Giorgi
approach. Another constructive proof by the Moser-Kru\v{z}kov
approach is proposed in~\cite{Guerand-Imbert}. The weak Harnack
inequality was obtained for the long-range Boltzmann equation
in~\cite{MR4049224}, and was proved for the kinetic
Fokker-Planck equations considered in this paper
in~\cite{Guerand-Imbert} by the Moser-Kru\v{z}kov approach.
\item As compared to that in~\cite{Guerand-Imbert}, our approach
is based on trajectorial arguments and does not require working
on the logarithm of the solution or the so-called inkspot
lemma. Our Poincaré inequality (Proposition~\ref{p:HPI
lemma}) and measure-to-pointwise estimate (Lemma
\ref{l:increase}) take into account a gap in time which removes
the requirement for the sub-solution to be considered in a
large domain. Our Poincaré inequality also holds without an
information in measure around the center of the cylinder as
in~\cite{Guerand-Imbert}.
\end{enumerate}
\end{remarks}

\begin{theorem}[Hölder continuity]
\label{t:holder}
There is $\alpha \in (0,1)$, computable from the proof and only
depending on $\lambda, \Lambda$ and $\|S\|_{L^{\infty}}$, such
that any weak solution $f$ of \eqref{e:main}--\eqref{e:hyp-coef}
in~$Q_2$ satisfies\vspace*{-3pt}
\begin{align*}
[f]_{C^\alpha (Q_{1})}
& := \sup_{\substack{z_1,z_2 \in Q_1\\ z_1 \not =
z_2}} \frac{|f(z_1) - f(z_2)|}{|z_1-z_2|^{\alpha}} \\
& \lesssim_{\lambda,\Lambda}
\bigl(1+\|S\|_{L^\infty(Q_2)} \bigr)
\bigl( \|f\|_{L^2 (Q_2)} + \|S\|_{L^\infty(Q_2)} \bigr).
\end{align*}
\end{theorem}

\begin{remark}
The Hölder continuity was first proved in \cite{wz09,wz11}
(with constructive method) and this proof is revisited and
simplified in~\cite{Guerand-Imbert}, including ideas and
methods from~\cite{moser,kru63,kru64}. An alternative
non-constructive proof was proposed in~\cite{gimv} following
the De~Giorgi method \cite{DG56}: the non-constructive part was
the intermediate-value lemma and we provide here a new
constructive argument.
\end{remark}

\subsection{Structure of the method}

The core of our proof is, given $f$ sub-solution
to~\eqref{e:main}--\eqref{e:hyp-coef} with $S=0$:\vspace*{-5pt}\enlargethispage{.5\baselineskip}%
\begin{align*}
f \in L^\zeta, \ \zeta >0
& \quad \xrightarrow{(1)}
\quad f \in L^\infty \cap L^1_{t,v}W^{\sfrac13-0,1}_x \\[-5pt]
& \quad \xrightarrow{(2)} \quad
\text{Weak Poincaré inequality in $L^1$} \\[-5pt]
& \quad \xrightarrow{(3)}
\quad \text{Intermediate-Value Lemma (Theorem~\ref{t:IVL}})\\[-5pt]
& \quad \xrightarrow{(4)}
\quad \text{Measure-to-pointwise estimate} \\[-5pt]
& \quad \xrightarrow{(5)} \quad \text{Weak log-Harnack
estimate} \\[-5pt]
& \quad \xrightarrow{(6)}
\quad \text{Weak Harnack estimate.}
\end{align*}
Once these steps are proved, it is immediate to prove the Harnack
inequality for solutions by combining the weak Harnack inequality
for super-solutions and step~(1) for sub-solutions. The Hölder
continuity follows classically (see Subsection~\ref{ss:Holder})
from either the measure-to-pointwise estimate applied to both
sub-solutions $f$ and $-f$, or from the Harnack inequality.
Step~(1) (Section~\ref{sec:int}) is semi-novel: it elaborates
upon ideas in~\cite{pp} to prove the first Lemma of De~Giorgi as
well as a gain of Sobolev regularity with the help of Kolmogorov
fundamental solutions. Step~(2) (Proposition~\ref{p:HPI lemma})
is the most novel step and introduces an argument based on
trajectories and the previous Sobolev regularity to ``noise'' the
$x$-dependency of the trajectories. Step~(3) (proof in
Subsection~\ref{ss:proof-ivl}) is novel and based on simple
energy estimates. Step~(4) (Lemma~\ref{l:increase} in
Subsection~\ref{ss:m2p}) is standard and sketched for the sake of
obtaining quantitative constants. Step~(5) (in
Section~\ref{sec:Harnack}) is semi-novel but immediate when
constants are quantified properly in the previous steps. Step~(6)
(in Section~\ref{sec:Harnack}) is novel in the context of
hypoelliptic equations but inspired from elliptic
equations~\cite{MR3565366}; it uses an induction, Vitali's
covering lemma and Step~(5) at every scale.

\subsubsection*{Acknowledgements} The authors are grateful to
C.\,Imbert for the inspirational interactions, as well as for
specific help with the literature and the comparison with the
Moser-Kru\v{z}kov approach in \cite{Guerand-Imbert}. The second
author would also like to thank L.\,Silvestre who pointed out
several years ago how Kolmogorov fundamental solutions were used
in~\cite{pp} to replace averaging lemma, which was the starting
point of our Section~\ref{sec:int} (and is also used
in~\cite{MR4049224}).


\section{Integral estimates revisited}
\label{sec:int}

In this section, we briefly revisit estimates from~\cite{pp,gimv}
on the gain of inte\-grability for sub-solutions (the kinetic
counterpart to the first lemma of De~Giorgi) and the low-order
Sobolev regularity estimate for sub-solutions, first mentioned
in~\cite{gimv}. We provide new proofs based on fundamental
solutions which, albeit variants of existing ones, seem simpler
and optimal.

\Subsection{The energy estimate}

\begin{proposition}[Energy estimate]
\label{prop:EE}
Let $f$ be a non-negative weak sub-solution
to~\eqref{e:main}--\eqref{e:hyp-coef} in an open set
$\cU \in \R^{1+2d}$. Given any
$Q_r(z_0) \subset Q_R(z_0) \subset \cU$ with $0<r<R$, one has
\begin{align*}
\int_{Q_r(z_0)} |\nabla_v f|^2
\lesssim_{\lambda,\Lambda} \cC(r,R,v_0) \biggl(
\int_{Q_R(z_0)} f^2 + \int_{Q_R(z_0)} f |S| \biggr),
\end{align*}
where $z_0=(t_0,x_0,v_0)$,
$Q_r^\tau (z_0) = \{(x,v)\in\R^{2d} :(\tau,x,v)\in
Q_r(z_0) \} $, and
\begin{equation}
\label{eq:C}
\cC(r,R,v_0) := \Bigl( 1 + \frac{1}{(R-r)^2} +
\frac{|v_0|+R}{(R-r)r^2} +\frac{1}{(R-r)r} \Bigr).
\end{equation}
\end{proposition}

\begin{proof}[Proof of Proposition~\ref{prop:EE}]
Consider $\varphi$ a smooth function valued in $[0,1]$ that is
equal to $1$ on $Q_r(z_0)$ and $0$ outside $Q_R(z_0)$. In order
to use $f\varphi^2$ as a test function, we argue by
density. Introduce
\begin{equation*}
\psi_n * [f \varphi] \varphi (z) :=
\int_{z'\in\R^{2d+1}} \psi_n \left(t-t',x-tv-(x'-t'v'),v-v'
\right) f(z') \varphi(z') \varphi(z),
\end{equation*}
where $\psi_n(t,x,v)=n^{4d+2} \psi(n^2 t,n^3 x, n v)$ and
$\psi(t,x,v):= \pi^{-d-\sfrac{1}{2}} e^{-t^2-|x|^2-|v|^2}$.
Then
\begin{align*}
I_n:= \big\langle \cT f, \psi_n * [f \varphi] \varphi
\big\rangle
& = \big\langle f, \psi_n * [(\cT f) \varphi] \varphi
\big\rangle \\
& = \frac{1}{2} \left[\big\langle \cT f, \psi_n * [f \varphi] \varphi
\big\rangle + \big\langle f, \psi_n * [(\cT f) \varphi] \varphi
\big\rangle \right] \\
& = \frac{1}{2}\, \bigl[-\big\langle f, \psi_n * [f \varphi] (\cT\varphi)
\big\rangle -\big\langle f, \psi_n * [f (\cT \varphi)] \varphi
\big\rangle \bigr],
\end{align*}
which converges to
$-\frac{1}{2} \big\langle f^2, \cT\varphi^2 \big\rangle$ as $n
\to \infty$. The
other terms in the inequation converge thanks to the bound
$f \in L^\infty ((a,b);L^2 (\Omega_x \times \Omega_v)) \cap L^2
((a,b) \times \Omega_x;H^1(\Omega_v))$. We deduce
\begin{align*}
\lambda \int_{Q_R(z_0)} |\nabla_v f|^2
\varphi^2 \dd z
& \le \int_{Q_R(z_0)} f^2 \Big( |\partial_t \varphi|
\varphi + (|v_0|+R) | \nabla_x \varphi| \varphi \Big) \dd z\\
& \qquad + \Lambda \int_{Q_R(z_0)} |\nabla_v
f| |\nabla_v \varphi| f \varphi \dd z \\
& \qquad + \Lambda
\int_{Q_R(z_0)} |\nabla_v f| f \varphi^2 \dd z +
\int_{Q_R(z_0)} f |S| \varphi^2 \dd z.
\end{align*}
The result follows from Cauchy-Schwarz' inequality
and
\[
|\partial_t \varphi| \lesssim \frac{1}{(R-r)r},\quad
|\nabla_x \varphi| \lesssim \frac{1}{(R-r)r^2},\quad
|\nabla_v \varphi| \lesssim \frac{1}{(R-r)}.\qedhere
\]
\end{proof}

\subsection{Integral estimates on Kolmogorov fundamental
solutions}

We denote $\cK := \cT - \Delta_v$.
\begin{lemma}[Estimates on the fundamental solution with constant
coefficients]
\label{lem:Kolmogorov}
Consider $f \ge 0$ locally integrable so that
$\cK f = \nabla_v \cdot F_1 + F_2 - m$ with
$F_1,F_2 \in L^1 \cap L^2(\R_- \times \R^{2d})$ and
$0 \le m \in M^1(\R_- \times \R^{2d})$ (a non-negative measure
with finite mass) and where $F_1, F_2$ and $m$ have compact
support in time included in some $\left(-\tau,0\right]$. Then
there for any $p \in [2,2+\sfrac1d)$ and
$\sigma \in [0,\sfrac13)$
\begin{equation}
\label{eq:Kolm1}
\| f\|_{L^{p}(\R_- \times \R^{2d})}\lesssim_{\tau,\lambda,\Lambda}
\left(2+\Psfrac1d-p\right)^{-1}
\bigl[ \left\| F_1 \right\|_{L^2(\R_- \times \R^{2d})} +
\left\| F_2 \right\|_{L^2(\R_- \times \R^{2d})} \bigr],
\end{equation}
\begin{multline}\label{eq:Kolm2}
\| f\|_{L^1_{t,v} W^{\sigma,1}_x(\R_- \times \R^{2d})} \\
\lesssim_{\tau,\lambda,\Lambda}
\left(\Psfrac13-\sigma \right)^{-1} \Bigl[ \|
F_1 \|_{L^1(\R_- \times \R^{2d})} + \left\| F_2
\right\|_{L^1(\R_- \times \R^{2d})} + \| m \|_{M^1(\R_- \times
\R^{2d})} \Bigr].
\end{multline}
\end{lemma}

\begin{proof}[Proof of Lemma~\ref{lem:Kolmogorov}]
We use the fundamental solution computed by Kolmogorov
in~\cite{kolmogorov} (see for
instance~\cite[App.\,A]{MR4113786} for details):
\begin{align*}
\forall \, t \in \R_-, \ x,v &\in \R^d, \\
f(t,x,v) &\phantom{:}=\int_{(t',x',v')\in \R^{2d+1}} G(t-t', x-x'
-(t-t')v',v-v') (\cK f)(t',x',v'), \\
\forall \, t \ge 0, \ x,v &\in \R^d, \\
G(t,x,v) &:=
\begin{dcases}
\Bigl( \frac{3}{4\pi^2 t^4} \Bigr)^{\sfrac{d}{2}}
\exp \Bigl[ - \frac{3 \left| x - \Psfrac{t}{2} v
\right|^2}{t^3} - \frac{|v|^2}{4t} \Bigr] & \text{if } t>0,\\
0 & \text{if } t\leq 0.
\end{dcases}
\end{align*}
Since $f$ and $G$ are non-negative, we deduce that
\begin{multline*}
0 \le f(t,x,v) \le\int_{(t',x',v')\in \R^{2d+1}} G(t-t', x-x'-(t-t')v',v-v')\\
\cdot\bigl[(\nabla_{v'} \cdot F_1) (t',x',v')+ F_2 (t',x',v') \bigr]
\end{multline*}
and since
\begin{multline*}
\forall \, t \ge 0, \ x,v \in \R^d, \\[-5pt]
|\nabla_v G(t,x,v)| +t |\nabla_x G(t,x,v)|
\lesssim t^{-2d-\sfrac12}
\exp \Bigl[ - \frac{3 \left| x - \Psfrac{t}{2} v
\right|^2}{2t^3} - \frac{|v|^2}{8t} \Bigr]
\end{multline*}
we have
$\nabla_v G, t \nabla_x G \in L^{\frac{2d+1}{2d+1/2}-0}
(\left(0,\tau\right) \times \R^{2d})$ and therefore by
integration by parts
\begin{multline*}
f(t,x,v) \le \int_{(t',x',v')\in \R^{2d+1}}
\nabla_{v'} G(t-t',x-x'-(t-t')v',v-v') F_1(t',x',v')\\
+ \int_{(t',x',v') \in \R^{2d+1}}
G(t-t',x-x'-(t-t')v',v-v') F_2(t',x',v')
\end{multline*}
and Young's convolution inequality (which works in unimodular
spaces like $(\R^{2d+1},\circ)$ with the Lebesgue measure), we
deduce, by tracking down the dependency in $p$ of the constant
\begin{multline*}
\forall \, p \in \left[2,2+\sfrac1d\right), \\
\| f \|_{L^{p}(\R_- \times \R^{2d})}\lesssim_\tau \left( 2+ \Psfrac1d - p \right)^{-1} \bigl[ \|
F_1 \|_{L^2(\R_- \times \R^{2d})} + \left\| F_2
\right\|_{L^2(\R_- \times \R^{2d})} \bigr].
\end{multline*}
(The threshold $2+\sfrac1d$ is likely to be optimal.) This
proves~\eqref{eq:Kolm1}. To prove~\eqref{eq:Kolm2} split
\begin{align*}
G = G_\var + G_\var^\bot \quad \text{ with } \quad
G_\var(t,x,v) := \chi \left( \sfrac{t}{\var} \right) G(t,x,v),
\end{align*}
where $\var >0$ and $\chi$ is a smooth function on $\R_+$
valued in $[0,1]$ equal to $1$ in $[0,1]$ and~$0$ on
$[2,+\infty)$.
\Changel
We have the following simple estimates for every
$l \in \mathbb{N}$
\begin{align*}
\left|\nabla_x^l G_\var^\bot(t,x,v)\right| &\lesssim_l
\var^{-\Psfrac32 l}t^{-2d}
\exp \Bigl[ - \frac{3 \left| x - \psfrac{t}{2} v
\right|^2}{2t^3} - \frac{|v|^2}{8t} \Bigr] \\
 \left|\nabla_v \nabla_x^l G_\var^\bot(t,x,v)\right|
+ t\bigl|\nabla_x \nabla_x^l &G_\var^\bot(t,x,v) \bigr|
\\[-3pt]
&\lesssim_l \var^{-\Psfrac32 l -\sfrac12} t^{-2d}
\exp \Bigl[ - \frac{3 \left| x - \Psfrac{t}{2} v
\right|^2}{2t^3} - \frac{|v|^2}{8t} \Bigr]
\end{align*}
which straightforwardly implies (assuming $\tau \ge 1$ and
$\var <1$ wlog)
\begin{align*}
\left\| G_\var^\bot \right\|_{L^1_{t,v}((0,\tau)
\times \R^d; W^{l,1}_x(\R^d)))} +
\left\| \nabla_v G_\var^\bot \right\|_{L^1_{t,v}((0,\tau)
\times \R^d ; W^{l,1}_x(\R^d))}\hspace*{1mm}& \\
+ \left\|t\nabla_x G_\var^\bot \right\|_{L^1_{t,v}((0,\tau)
\times \R^d ; W^{l,1}_x(\R^d))}
&\lesssim_l \tau \var^{-\Psfrac32 l -\sfrac12} \\
\| G_\var \|_{L^1((0,\tau) \times \R^{2d})} + \left\|
\nabla_v G_\var \right\|_{L^1((0,\tau) \times \R^{2d})}
+ \left\|t\nabla_x G_\var \right\|_{L^1((0,\tau) \times \R^{2d})}
&\lesssim \tau \var^{\sfrac12}.
\end{align*}
The splitting $G=G_\var + G_\var^\bot$ yields
$f = f_\var + f_\var^\bot$, and the convolution inequality
$M^1 * L^1 \to L^1$ implies
\begin{align*}
& \| f_\var \|_{L^1(\R_- \times \R^{2d})} \lesssim
\tau \var^{\sfrac12} \bigl( \| F_1 \|_{L^1(\R_- \times \R^{2d})} +
\left\| F_2 \right\|_{L^1(\R_- \times \R^{2d})}
+ \| m \|_{M^1(\R_- \times
\R^{2d})} \bigr) \\
& \bigl\| f_\var^\bot \bigr\|_{L^1_{t,v}W^{l,1}_x(\R_-
\times \R^{2d})} \\
& \hphantom{\| f_\var \|_{L^1(\R_- \times \R^{2d})}} \lesssim
\tau \var^{-\Psfrac32 l -\sfrac12}
\bigl[ \| F_1 \|_{L^1(\R_- \times \R^{2d})} \!+\!
\left\| F_2 \right\|_{L^1(\R_- \times \R^{2d})}
\!+\! \| m \|_{M^1(\R_- \times \R^{2d})} \bigr].
\end{align*}
Since this decomposition holds for all $\var >0$, it implies by
standard interpolation the estimate~\eqref{eq:Kolm2} for any
$\sigma \in [0,\sfrac{1}{3})$ (again the exponent is likely to
be optimal but in any case our constant degenerates as
$\sigma \to \sfrac{1}{3}$). In order to be self-contained let us
a give a short proof. Given $\sigma \in [0,\sfrac13)$, we
Fourier-transform and decompose dyadically, defining
$\langle \xi \rangle:= (1+|\xi|^2)^{\sfrac12}$
\begin{equation}
\label{eq:decomp}
\begin{aligned}
(1-\Delta_x)&^{\sfrac{\sigma}{2}} f(t,x,v) \\
&= \int_{\xi,y \in \R^d} e^{i \xi
\cdot (x-y)} \langle \xi \rangle^\sigma
f(t,y,v)
= \sum_{k \ge -1} \int_{\xi,y \in \R^d} e^{i \xi
\cdot (x-y)} a_k(\xi) f(t,y,v) \\
&= \sum_{k \ge -1} \int_{\xi,y \in \R^d} e^{i \xi
\cdot (x-y)} B_k(\xi) (1-\Delta_y)^{\sfrac{l}{2}} f(t,y,v),
\end{aligned}
\end{equation}
where $a_k(\xi) := \langle \xi \rangle^\sigma \varphi_k$ and
$B_k(\xi) := \langle \xi \rangle^{\sigma-l} \varphi_k$, and
where we have defined in the standard way
$\varphi_k(\xi) := [ \chi(2^{-k}\xi) - \chi(2^{-k+1} \xi) ]$
for $k \ge 0$ with $\chi$ a smooth function valued in $[0,1]$
and equal to $1$ in $B(0,1)$ and $0$ outside $B(0,2)$, and
$\varphi_{-1}(\xi) =\nobreak \sum_{k \le -1} [ \chi(2^{-k}\xi) -\nobreak
\chi(2^{-k+1} \xi) ]$. For a given $F=F(y)$ one has
\begin{align*}
\int_{x \in \R^d} \left| \int_{\xi,y \in \R^d} e^{i \xi
\cdot (x-y)} a_k(\xi) F(y)
\right| &\lesssim 2^{k\sigma} \| F \|_{L^1}, \\
\int_{x \in \R^d} \left| \int_{\xi,y \in \R^d} e^{i \xi
\cdot (x-y)} B_k(\xi) (1-\Delta_y) F(y)
\right| &\lesssim 2^{k(\sigma-l)} \| F \|_{W^{l,1}}
\end{align*}
by splitting the integrand into $|x-y| \le 2^{-k}$ and
$|x-y| > 2^{-k}$ and integrating by parts the operator
$\Delta_\xi^{\sfrac{\ell}{2}}$ with $\ell$ even and strictly
greater than $d$. We then use the
decomposition~\eqref{eq:decomp} in the ``$a_k$'' form on
$f_\var$ and in the ``$B_k$'' form on $f_\var^\bot$, and with
a $\var=\var_k$ depending on $k$ defined below:
\begin{align*}
\bigl\| (1-\Delta_x)^{\sfrac{\sigma}{2}} f \bigr\|_{L^1(\R_-
\times \R^{2d})}
& \lesssim \tau \sum_{k \ge -1} \bigl( \var_k
^{\sfrac12} 2^{k\sigma} + \var_k^{-\Psfrac{3}{2}l -\sfrac12}
2^{k(\sigma-l)} \bigr) \\
&\hspace*{3mm}
\times \bigl[ \| F_1 \|_{L^1(\R_- \times \R^{2d})}
+ \left\| F_2 \right\|_{L^1(\R_- \times \R^{2d})}
+ \| m \|_{M^1(\R_- \times \R^{2d})} \bigr]\\
& \lesssim
\frac{\tau}{\delta} \bigl[ \| F_1 \|_{L^1(\R_- \times \R^{2d})}
+ \left\| F_2 \right\|_{L^1(\R_- \times \R^{2d})}
+ \| m \|_{M^1(\R_- \times \R^{2d})} \bigr]
\end{align*}
with the choice $\sigma = \sfrac13 - \delta \in [0,\sfrac13)$ and
$\var_k := 2^{-2k(\sfrac13-\sfrac{\delta}{2})}$ and
$l> 1+\sfrac{4}{9\delta}$. This concludes the proof.
\Changelback
\end{proof}

\subsection{Integral estimates for sub-solutions}

We combine the previous lemma with a localization argument and
the energy estimate to get the
\begin{proposition}[Integral regularization estimates for
non-negative sub-solutions]
\label{prop:Gain sol}
Let $f$ be a non-negative weak sub-solution
to~\eqref{e:main}--\eqref{e:hyp-coef} in an open set
$\cU \in \R^{1+2d}$. Given any
$Q_r(z_0) \subset Q_R(z_0) \subset \cU$ with $0<r<R \le 1$, and
any $p \in [2,2+\sfrac1d)$ and $\sigma \in [0,\sfrac13)$, $f$
satisfies
\begin{align}
\label{eq:gain-int}
\|f\|_{L^{p}(Q_r(z_0))} \lesssim ( 2+ \sfrac1d - p)^{-1}
\cC'(r,R,v_0) \bigl[ \norm{f}_{L^{2}(Q_R(z_0))}
+ \|S\|_{L^2(Q_R(z_0))} \bigr] ,\\
\label{eq:gain-reg}
\| f \|_{L^1_{t,v}W^{\sigma,1}_x(Q_r(z_0))} \lesssim \left(
\sfrac13 - \sigma \right)^{-1}
\cC''(r,R,v_0) \bigl[ \norm{f}_{L^{2}(Q_R(z_0))}
+ \|S\|_{L^2(Q_R(z_0))} \bigr],
\end{align}
where $\cC$ was defined in~\eqref{eq:C} and
\begin{align*}
\cC'(r,R,v_0) &:= \Bigl( 1+ \frac{1}{R-r} \Bigr)
\cC(r,R,v_0),\\
\cC''(r,R,v_0) &:= R^{1+2d} \Bigl( 1+ \frac{1}{R-r} \Bigr)
\cC(r,R,v_0).
\end{align*}
\end{proposition}

\begin{proof}[Proof of Proposition~\ref{prop:Gain sol}]
Since $f$ is a sub-solution to~\eqref{e:main}, there is a
non-negative measure $\bar m \ge 0$ so that
\begin{align*}
\cT f =\nabla_v\cdot (A\nabla_v f) + B \cdot \nabla_v f
+ S - \bar m.
\end{align*}
Consider $\varphi_1$ smooth valued in $[0,1]$ and equal to $1$
on $Q_r(z_0)$ and $0$ outside $Q_{r+\psfrac{R-r}{2}}(z_0)$ and
$g_1:=\varphi_1 f$. The latter satisfies
\begin{equation}
\label{eq:ineqKg}
\begin{aligned}
\cK g_1 &= \nabla_v \cdot F_1+F_2 - m \\
&\text{with }
\begin{dcases}
m:= \bar m \varphi_1, \\
F_1 := (A\nabla_v f) \varphi_1-(\nabla_v
f) \varphi_1 - f \nabla_v \varphi_1, \\
F_2:= -A\nabla_v f\cdot \nabla_v \varphi_1+ \left( B
\cdot \nabla_v f \right)\varphi_1 + S\varphi_1 + f \cT
\varphi_1.
\end{dcases}
\end{aligned}
\end{equation}
The energy estimate in Proposition \ref{prop:EE} implies
\begin{align*}
& \| F_1 \|_{L^2(\R_- \times \R^{2d})} + \| F_2 \|_{L^2(\R_-
\times \R^{2d})} \\
& \hspace{1cm} \lesssim \Bigl( 1+\frac{1}{R-r} \Bigr)
\cC\Bigl(r+\frac{R-r}{2}, R,v_0\Bigr)
\left( \| f \|_{L^2\left(Q_R(z_0)\right)}
+ \| S \|_{L^2(Q_R(z_0))} \right) \\
& \hspace{1cm} \lesssim \cC'(r,R,z_0) \bigl( \| f
\|_{L^2(Q_R(z_0))} + \| S \|_{L^2(Q_R(z_0))} \bigr),
\end{align*}
which, combined with~\eqref{eq:Kolm1},
shows~\eqref{eq:gain-int}.

Consider then $\varphi_2$ smooth valued in $[0,1]$ and equal to
$1$ on $Q_{r+\psfrac{R-r}{2}}(z_0)$ and~$0$ outside $Q_R(z_0)$
and $g_2:=\varphi_2 f$. The function $g_2$ satisfies a similar
equation as $g_1$ in~\eqref{eq:ineqKg}, with $\varphi_2$
replacing $\varphi_1$. Integrating this equation simply against
$1$ yields (thanks to the cancellation of divergence terms)
\begin{align*}
\| \bar m \|_{M^1\left(Q_{r+\psfrac{R-r}{2}}(z_0)\right)}
& \lesssim
\| \varphi_2 m \|_{M^1(\R_- \times \R^{2d})} \\
& \lesssim \int_{Q_{r+\psfrac{R-r}{2}}(z_0)}
\hspace*{-2mm}[ -A\nabla_v
f\!\cdot\! \nabla_v \varphi_2 + \left( B \!\cdot\! \nabla_v f
\right)\varphi_2 + S\varphi_2 + f \cT \varphi_2] \\
& \lesssim \cC\Bigl(r+\frac{R-r}{2},R,v_0 \Bigr) \| f
\|_{L^1(Q_R(z_0))} + \| S \|_{L^1(Q_R(z_0))}\\
& \lesssim \cC\left(r,R,v_0
\right) \bigl[ \| f \|_{L^2(Q_R(z_0))} + \| S
\|_{L^2(Q_R(z_0))} \bigr].
\end{align*}
Combined with~\eqref{eq:Kolm2} and (thanks to the localization)
\begin{align*}
\| F_1 \|_{L^1(\R_- \times \R^{2d})} + \| F_2 \|_{L^1(\R_-
\times \R^{2d})} \lesssim \| F_1 \|_{L^2(\R_- \times
\R^{2d})} + \| F_2 \|_{L^2(\R_- \times \R^{2d})},
\end{align*}
it implies~\eqref{eq:gain-reg}.
\end{proof}

\subsection{Iterated gain of integrability for sub-solutions}

We give a short proof of this result first obtained in
\cite[Th.\,1.2]{pp} and then proved
differently~\cite[Th.\,12]{gimv}. This is the counterpart of
the ``first lemma of De~Giorgi'' for elliptic equations, in the
context of kinetic hypoelliptic equations. We allow for an
initial integrability~$L^\zeta$ with exponent $\zeta \in (0,2)$
(such extension is well-known for elliptic equations).

\begin{proposition}[Upper bound for sub-solutions]
\label{prop:1st lemma}
Let $f$ be a non-negative weak sub-solution
to~\eqref{e:main}--\eqref{e:hyp-coef} in an open set
$\cU \in \R^{1+2d}$. Given any
$Q_r(z_0) \subset Q_R(z_0) \subset \cU$ with $0<r<R \le 1$, and
$\zeta >0$, $f$ satisfies
\begin{equation*}
\| f \|_{L^\infty(Q_r(z_0))} \lesssim_{\lambda,\Lambda}
\Bigl( \frac{1+|v_0|}{r^2(R-r)^3}\Bigr)^{\psfrac{1+4d}{\zeta}}
\bigl[ \|f\|_{L^\zeta(Q_R(z_0))} + \|S\|_{L^\infty(Q_R(z_0))} \bigr].
\end{equation*}
\end{proposition}

\begin{proof}[Proof of Proposition~\ref{prop:1st lemma}]
Fix $p_0:=2+\sfrac{1}{2d}$ and define $q:=\sfrac{p_0}{2}$ and
$q_n := q^n$. Consider $\beta_{n,k}$ on $\R_+$ with
$\beta_{n,k}' \ge 0$ and $\beta''_{n,k} \ge 0$ both bounded and
so that $\beta_{n,k}(z) \to\nobreak z^{q_n}$ as $k \to \infty$ and
$\beta_{n,k}(z) \lesssim z^{q_n}$ and
$\beta_{n,k}'(z) \lesssim z^{q_n-1}$ uniformly in $k \in
\N^*$. Definition~\ref{d:weak} implies that
$\beta_{n,k}(f)$ is a weak sub-solution with source term
$S_{n,k} := \beta_{n,k}'(f) S$. Define $r_0=R$ and
$r_n := r_{n-1} - \delta n^{-2}$ with
$\delta = \frac12( \sum_{k \ge 1} k^{-2}
)^{-1}(R-r)$. Since $p_0 \in [2,2+\sfrac1d)$, the
estimate~\eqref{eq:gain-int} implies for all $n \ge 1$
\begin{align*}
\|\beta_{n,k}(f)\|_{L^{p_0}(Q_{r_n}(z_0))}
& \lesssim
\cC'(r_n,r_{n-1},v_0) \bigl[
\|\beta_{n,k}(f)\|_{L^{2}(Q_{r_{n-1}}(z_0))}
+ \|S_{n,k}\|_{L^2(Q_{r_{n-1}}(z_0))} \bigr] \\
& \lesssim \frac{(1+|v_0|)n^6}{r^2(R-r)^3}
\bigl[ \|\beta_{n,k}(f)\|_{L^{2}(Q_{r_{n-1}}(z_0))}
+ \|S_{n,k}\|_{L^2(Q_{r_{n-1}}(z_0))} \bigr]
\end{align*}
for $n \ge 1$, which means by taking $k \to \infty$ and coming
back to $f$
\begin{multline*}
\|f\|_{L^{2q_{n+1}}(Q_{r_n}(z_0))} \\
\lesssim \Bigl(
\frac{(1+|v_0|)n^6}{r^2(R-r)^3} \Bigr)^{\sfrac{1}{q^n}}
2^{-1+\sfrac{1}{q_n}} \Bigl[
\|f\|_{L^{2q_n}(Q_{r_{n-1}}(z_0))} + \|
f\|_{L^{2q_n}(Q_{r_{n-1}}(z_0))}^{1-\sfrac{1}{q_n}}
\|S\|_{L^\infty(Q_R(z_0))}^{\sfrac{1}{q_n}} \Bigr] \\
\lesssim \Bigl(
\frac{(1+|v_0|)n^6}{r^2(R-r)^3} \Bigr)^{\sfrac{1}{q^n}}
\Bigl[ \Bigl( 1 + \frac{1}{q_n} \Bigr)
\|f\|_{L^{2q_n}(Q_{r_{n-1}}(z_0))} + \frac{1}{q_n}
\|S\|_{L^\infty(Q_R(z_0))} \Bigr],
\end{multline*}
assuming by induction
$\|f\|_{L^{2q_n}(Q_{r_{n-1}}(z_0))} < +\infty$. The convergence
of the infinite product then implies
\begin{equation*}
\|f\|_{L^\infty(Q_r(z_0))}
\lesssim \Bigl( \frac{1+|v_0|}{r^2(R-r)^3}
\Bigr)^{1+4d} \bigl[ \|f\|_{L^2(Q_R(z_0))} +
\|S\|_{L^\infty(Q_R(z_0))} \bigr].
\end{equation*}
This proves the claim when $\zeta \ge 2$. To prove it when
$\zeta \in (0,2)$, we deduce from the previous estimate
\begin{multline*}
\|f\|_{L^\infty(Q_r(z_0))} + \|S\|_{L^\infty(Q_r(z_0))}
\\
\lesssim \Bigl( \frac{1+|v_0|}{r^2(R-r)^3}
\Bigr)^{1+4d} \bigl[ \|f\|_{L^\infty(Q_R(z_0))}^{1-\zeta}
\|f\|_{L^\zeta(Q_R(z_0))}^\zeta +
\|S\|_{L^\infty(Q_R(z_0))}\bigr]
\end{multline*}
and thus by Young inequality, the quantity $A(r) :=
\|f\|_{L^\infty(Q_r(z_0))} + \|S\|_{L^\infty(Q_r(z_0))}$ satisfies,
for some $C>0$,
\begin{align*}
A(r) \le \frac12 A(R) + C \Bigl( \frac{1+|v_0|}{r^2(R-r)^3}
\Bigr)^{\psfrac{1+4d}{\zeta}} \bigl[ \|f\|_{L^\zeta(Q_R(z_0))} +
\|S\|_{L^\infty(Q_R(z_0))}\bigr].
\end{align*}
Introducing an (increasing this time) sequence of radii
$r_n := r_{n-1} + \delta n^{-2}$ we obtain by induction
\begin{align*}
A(r_n) &\le \Psfrac12 A(r_{n+1})\\
& + C n^{\psfrac{2+8d}{\zeta}}
\Bigl( \frac{1+|v_0|}{r^2(R-r)^3}
\Bigr)^{\psfrac{1+4d}{\zeta}} \bigl[ \|f\|_{L^\zeta(Q_R(z_0))} +
\|S\|_{L^\infty(Q_R(z_0))} \bigr], \\
A(r_0) &\le \left( \sfrac12 \right)^n A(r_{n+1}) \\
& \hspace{2mm}+
C \biggl( \sum_{k=1}^n \frac{k^{\psfrac{2+8d}{\zeta}}}{2^k} \biggr)
\Bigl( \frac{1+|v_0|}{r^2(R-r)^3}
\Bigr)^{\psfrac{1+4d}{\zeta}} \bigl[ \|f\|_{L^\zeta(Q_R(z_0))} +
\|S\|_{L^\infty(Q_R(z_0))}\bigr],
\end{align*}
which yields the result by taking $n \to \infty$ in the right
hand side.
\end{proof}

\section{Intermediate-Value Lemma and oscillations}

\subsection{Weak Poincaré inequality}

The adjective `weak' refers to the small additional $L^2$ error
term below.

\begin{proposition}[Hypoelliptic Poincaré inequality with error]
\label{p:HPI lemma}
Given any $\var\in (0,1)$ and $\sigma \in (0,\sfrac13)$, any
non-negative sub-solution $f$ to
\eqref{e:main}--\eqref{e:hyp-coef} on $Q_5$ satisfies
\begin{multline}
\label{HPI inequality}
\bnorm{\bigl(f-\langle f
\rangle_{Q_1^{-}} \bigr)_+}_{L^{1}(Q_1)}\\
\lesssim_{\lambda,\Lambda}
\frac{1}{\varepsilon^{d+2}} \norm{\nabla_v f}_{L^1 (Q_5)}
+ \varepsilon^\sigma \left(
\sfrac13 - \sigma \right)^{-1} \norm{f}_{L^2(Q_5)} +
\nnorm{S}_{L^2(Q_5)},
\end{multline}
where $Q_1^- := Q_1(-2,0,0) = (-3,-2] \times B_1 \times B_1$
and
$\langle f \rangle_{Q_1^{-}} := \fint_{Q_1^-} f :=
\frac{1}{|Q_1^-|} \int_{Q_1^-} f$.
\end{proposition}

\begin{remark}
The motivation for the following argument was \cite[Lem.\,10,
p.\,11]{vasseur}, where a simple quantitative proof of the
intermediate value lemma of De~Giorgi (also sometimes called De
Giorgi's isoperimetric inequality) is sketched in the elliptic
case, based on introducing the trajectory between two points of
the domain and using the vector field $\nabla_v$ to connect
them. We have to deal here with the hypoelliptic structure.
\end{remark}

\begin{proof}
Consider, for $\var \in (0,1)$, a smooth function
$\varphi_\var=\varphi_\var(y,w)$ which satisfies
$0\leq \varphi_\var \leq 1$ and has compact support in
$B_1^{2}$ and such that $\varphi_\var=1$ in
$B_{(1-\var)}\times B_{(1-\var)}$ and with
$|\nabla_y \varphi_\var|\lesssim \var^{-1}$ and
$|\nabla_w \varphi_\var|\lesssim \var^{-1}$. We then split the
integral to be estimated as follows
\begin{multline*}
\bnorm{\bigl(f-\langle f
\rangle_{Q_1^{-}}\bigr)_+}_{L^{1}(Q_1)} \lesssim
\bnorm{\bigl(f-\langle f \varphi_\var
\rangle_{Q_1^{-}} \bigr)_+}_{L^{1}(Q_1)}\\
\hspace*{-1.13cm}\lesssim
\int_{(t,x,v) \in Q_1}
\biggl\{ \fint_{(s,y,w) \in Q_1^-} \left[
f(t,x,v)-f(s,y,w) \right] \varphi_\var(y,
w) \biggr\}_{\!+} \\
\shoveright{+ \norm{f}_{L^1(Q_1)}
\fint_{Q_{1}^{-}}\left( 1 - \varphi_\var(y,w) \right)}\\
\lesssim \int_{(t,x,v) \in Q_1}
\biggl\{ \fint_{(s,y,w) \in Q_1^-} \left[
f(t,x,v)-f(s,y,w) \right] \varphi_\var(y,
w) \biggr\}_{\!+} + \var^{2d} \nnorm{f}_{L^2(Q_1)},
\end{multline*}
where we have used
$\langle f \varphi_\var \rangle_{Q_1^{-}}\leq \langle f
\rangle_{Q_1^{-}}$ and the Cauchy-Schwarz inequality.

Let us estimate the first term of the previous
inequality. Given $t,x,v$ fixed, we~decompose the trajectory
$(t,x,v)\to (s,y,w)$ into four sub-trajectories in
$Q_5$: a trajectory of length $O(\var)$ along $\nabla_x$, two
trajectories of length $O(1)$ along $\nabla_v$, and finally one
trajectory of length $O(1)$ along
$\cT := \partial_t + v \cdot \nabla_x$:
\begin{multline*}
(t,x,v)\underset{\nabla_x}{\to}
(t,x+\varepsilon w,v) \underset{\nabla_v}{\to}
\Bigl(t,x+\varepsilon w, \frac{x+\varepsilon w
-y}{t-s} \Bigr) \\
\underset{\cT}{\to} \Bigl(s,
y,\frac{x+\varepsilon w -y}{t-s} \Bigr)
\underset{\nabla_v}{\to} (s,y,w).
\end{multline*}
The first sub-trajectory is estimated by the {\it integral}
regularity $L^1_{t,v} W^{\sigma,1}_{x}$ proved
in~\eqref{eq:gain-reg}. The other trajectories are estimated
directly by the vector fields in the equation. The position
$x+\varepsilon w \in Q_2$ since $x,w \in B_1$ and
$\var \in (0,1)$. The velocity
$\pspfrac{x+\varepsilon w -y}{t-s} \in Q_3$ since $x,w,y \in B_1$
and $t-s\geq 1$ due to the definitions of~$Q_{1}^+$ and
$Q_{1}^{-}$, and this velocity yields a transport line from
$(t,x+\varepsilon w)$ to $(s,y)$. Note that we are implicitly
using the Hörmander commutator condition:
$\nabla_v,\cT, [\nabla_v,\cT]$ span all the vector fields on
$\R^{2d+1}$.

Decompose along the previous trajectories
\begin{align*}
f(t&,x,v)-f(s,y,w) \\
& = \bigl[ f(t,x,v)-f(t,x+\varepsilon w,v) \bigr] +
\Bigl[ f(t,x+\varepsilon
w,v)- f \Bigl(t,x+\varepsilon w,\frac{x+\varepsilon w -y}{t-s}
\Bigr) \Bigr] \\
& \qquad + \Bigl[ f \Bigl(t,x+\varepsilon w,\frac{x+\varepsilon w
-y}{t-s} \Bigr)- f \Bigl(s,y,\frac{x+\varepsilon w
-y}{t-s} \Bigr) \Bigr] \\
& \qquad + \Bigl[ f \Bigl(s,y,\frac{x+\varepsilon w
-y}{t-s} \Bigr)-f(s,y,w) \Bigr]
\end{align*}
and integrate against $\varphi_\var(y,w)$ on
$(s,y,w)\in Q_{1}^-$, which gives the four terms
\begin{align*}
I_1(t,x,v)
& := \int_{(s,y,w)\in Q_1^{-}}
\bigl[ f(t,x,v)-f(t,x+\varepsilon
w,v)\bigr] \varphi_\var(y,w),\\
I_2(t,x,v)
& := \int_{(s,y,w)\in Q_1^{-}}
\Bigl[ f(t,x+\varepsilon w,v)-f\Bigl(t,x+\varepsilon
w,\frac{x+\varepsilon w -y}{t-s}\Bigr)\Bigr]
\varphi_\var(y,w), \\
I_3(t,x,v)
& := \int_{(s,y,w)\in Q_1^{-}} \Bigl[ f\Bigl(t,x+\varepsilon
w,\frac{x+\varepsilon w
-y}{t-s}\Bigr)-f\Bigl(s,y,\frac{x+\varepsilon w
-y}{t-s}\Bigr)\Bigr] \varphi_\var(y,w), \\
I_4(t,x,v)
& := \int_{(s,y,w)\in Q_1^{-}}
\Bigl[ f\Bigl(s,y,\frac{x+\varepsilon w
-y}{t-s}\Bigr)-f(s,y,w)\Bigr] \varphi_\var(y,w).
\end{align*}

Regarding the term $I_2$, we use Taylor's formula and
$0\leq \varphi_\var\leq 1$ to deduce
\begin{align*}
I_2(t,x,v)
&\leq \int_{(s,y,w) \in Q_1^{-}} \int_{\tau\in [0,1]}
\Bigl(v-\frac{x+\varepsilon w-y}{t-s}\Bigr) \cdot \\
& \hspace{3.3cm} \times \nabla_v
f \Bigl(t,x+\varepsilon w, \tau v +
(1-\tau)\frac{x+\varepsilon w-y}{t-s} \Bigr)
\varphi_\var(y,w) \\
&\lesssim \int_{(s,y,w) \in Q_1^{-}} \int_{\tau\in [0,1]}
|\nabla_v f| \Bigl(t,x+\varepsilon w, \tau v +
(1-\tau)\frac{x+\varepsilon w-y}{t-s} \Bigr).
\end{align*}
Integrate then on $(t,x,v)\in Q_{1}^+$ to get
\begin{multline}
\label{estim I2}
\fint_{(t,x,v) \in Q_1} I_2 \\
\hspace*{-3mm}\lesssim \int_{(t,X,v)\in (-1,0)\times B_2 \times
B_1} \int_{(s,Y,w)\in (-3,-2)\times B_4 \times B_1}
\int_{\tau\in (0,1)} |\nabla_v f| \left(t,X, v+ (1-\tau) Y
\right) \\
\lesssim \int_{(t,X,V)\in (-1,0)\times B_2 \times
B_5} \int_{(s,Y,w)\in (-3,-2)\times B_4 \times B_1}
\int_{\tau\in (0,1)} |\nabla_v f| \left(t,X, V\right)
\lesssim \int_{Q_5} |\nabla_v f|,
\end{multline}
where we have used successively the following changes of
variables with bounded Jacobians:
\begin{equation*}
x\to X=x+\varepsilon w \in B_2, \quad
y\to Y= \frac{X -y}{t-s} -v\in B_4, \quad
v\to V=v+(1-\tau)Y\in B_5.
\end{equation*}

The term $I_4$ is treated like $I_2$:
\begin{align}
\label{estim I4}
\fint_{(t,x,v) \in Q_1} I_4 \lesssim \int_{Q_5} |\nabla_v f|.
\end{align}

Regarding the term $I_1 $, we perform the change of variable
$w \in B_1 \to x'=x+\varepsilon w \in B_\var(x)$ with
Jacobian $\var^{-d}$ and use the $L^1_{t,v} W^{\sigma,1}_{x}$
regularity of non-negative sub-solutions proved
in~\eqref{eq:gain-reg}:\vspace*{-3pt}
\begin{equation}
\label{estim I1}
\begin{aligned}
\fint_{(t,x,v)\in Q_1} I_1
& \lesssim \int_{(t,x,v)\in Q_1, \, (s,y,w)\in Q_1^{-}}
|f(t,x,v)-f(t,x+\varepsilon w,v) | \\
&\lesssim \int_{(t,x,v)\in Q_1, \, w \in B_1}
\frac{|f(t,x,v)-f(t,x+\varepsilon w,v)|}{|\varepsilon w
|^{d+\sigma}}\,|\varepsilon w|^{d+\sigma} \\
&\lesssim \var^\sigma \int_{(t,x,v)\in Q_1, \, x' \in B_2 }
\frac{|f(t,x,v)-f(t,x',v)|}{|x-x'|^{d+\sigma}} \\
& \lesssim \var^\sigma \norm{f}_{L^1_{t,v}W^{\sigma,1}_x(Q_2)}
\lesssim \var^\sigma \left( \sfrac13 -\sigma \right)^{-1}
\bigl[ \norm{f}_{L^2(Q_3)} + \norm{S}_{L^2(Q_3)} \bigr].
\end{aligned}
\end{equation}

Regarding the term $I_3$, we note first that
$\cT f \in L^2_{t,x} H^{-1}_v + M^1_{t,x,v}$ with finite norm
in $Q_R(z_0)$ (arguing as in proof of
Proposition~\ref{prop:Gain sol}). The Taylor formula between
$(t,x+\varepsilon w)$ and $(s,y)$ along $\cT$ thus holds in
weak form against $\varphi_\var$ thanks to the latter bounds
and the non-singular change of
variable~\eqref{eq:change-variables} discussed below:\vspace*{-5pt}
\begin{multline}\label{eq:Taylor-weak}
I_3(t,x,v)\\
= \int_{(s,y,w)\in Q_1^{-}} \Bigl[
f\Bigl(t,x+\varepsilon w,\frac{x+\varepsilon w
-y}{t-s}\Bigr)-f\Bigl(s,y,\frac{x+\varepsilon w
-y}{t-s}\Bigr)\Bigr] \varphi_\var(y,w)\\
\shoveleft{\hspace*{-1mm}\lesssim \int_{(s,y,w)\in Q_1^{-}} \int_{\tau \in
[0,1]} (t-s)}\\
\times \cT
f\Bigl(\tau t+(1-\tau)s,\tau (x+\varepsilon
w)+(1-\tau)y,\frac{x+\varepsilon w
-y}{t-s}\Bigr)\varphi_\var(y,w).
\end{multline}
We then use the fact that $f$ is a sub-solution
to~\eqref{e:main} in the distributional sense:\vspace*{-5pt}
\begin{align*}
I_3(t&,x,v)\\
&\lesssim \int_{(s,y,w)\in Q_1^{-}} \int_{\tau \in
[0,1]} (t-s) \\
&\hspace*{1.1cm}\times \nabla_v \cdot (A \nabla_v f)\Bigl(\tau
t+(1-\tau)s,\tau (x+\varepsilon
w)+(1-\tau)y,\frac{x+\varepsilon w
-y}{t-s}\Bigr)\varphi_\var(y,w)\\
&\hspace*{.2cm}+\int_{(s,y,w)\in Q_1^{-}} \int_{\tau \in
[0,1]} (t-s) \\
&\hspace*{1.1cm}\times
B\cdot \nabla_v f\Bigl(\tau t+(1-\tau)s,\tau
(x+\varepsilon w)+(1-\tau)y,\frac{x+\varepsilon w
-y}{t-s}\Bigr)\varphi_\var(y,w)\\
&\hspace*{.2cm}+ \int_{(s,y,w)\in Q_1^{-}} \int_{\tau \in
[0,1]} (t-s) \\
&\hspace*{1.1cm}\times S\Bigl(\tau t+(1-\tau)s,\tau (x+\varepsilon
w)+(1-\tau)y,\frac{x+\varepsilon w
-y}{t-s}\Bigr)\varphi_\var(y,w)\\
&:= I_{31}+I_{32}+I_{33}.
\end{align*}
Arguing as for $I_2$ and $I_4$, we have
\begin{align}
\label{estim I323}
\fint_{(t,x,v) \in Q_1} I_{32} \lesssim \int_{Q_5} \left|
\nabla_v f\right| \quad \text{ and } \quad
\fint_{(t,x,v) \in Q_1} I_{33} \lesssim \int_{Q_5} |S|,
\end{align}
where we performed consecutively the changes of variable
\begin{gather*}
y\to V=\frac{x+\varepsilon w-y}{t-s},\quad
x\to X= x+\varepsilon w -(1-\tau)(t-s)V,\\
s\to s'=t-s\quad \text{and}\quad t'\to t-(1-\tau)s'.
\end{gather*}
To estimate the remaining term $I_{31}$, we use the change of
variable
\begin{equation}
\label{eq:change-variables}
(y,w) \mto (Y,W) \quad \text{with }
Y:=\tau (x+\varepsilon w)+(1-\tau)y \text{ and } W:=\frac{x+\varepsilon w -y}{t-s}
\end{equation}
such that $(y,w)\mto (Y,W)$ is a bijection from the set
$(B_1)^2$ to the (diamond-shaped) set
\begin{equation*}
E:=E(\tau,\varepsilon,t,s,x) \subset B\left(\tau x, (1-\tau)
+ \tau \var\right)
\times B\Bigl(\frac{x}{t-s},\frac{1+\var}{t-s}\Bigr)
\subset B_2 \times B_3
\end{equation*}
with Jacobian $(\spfrac{\var}{t-s})^d$ and which maps
respective boundaries (to compute the Jacobian easily use the
formula
$\det \bigl( \begin{smallmatrix} A & B \\ C & D \end{smallmatrix}\bigr) =
\det(A - B D^{-1} C) \det D$). We deduce
\begin{multline*}
I_{31} = \frac{1}{\varepsilon^d} \int_{\tau\in
[0,1]}\int_{s\in (-3,-2), \, (Y,W)\in E}
(t-s)^{d+1}\nabla_v\cdot (A
\nabla_v f)\left(\tau
t+(1-\tau)s,Y,W\right)\\
\times \varphi_{\varepsilon} \Bigl(Y-\tau
(t-s)W,\frac{Y-x+(1-\tau)(t-s)W}{\varepsilon} \Bigr)
\end{multline*}
and we integrate by parts in $W$, using that $\varphi_\var=0$
on the boundary of $E(\tau,\var,t,s,x)$:
\begin{multline*}
I_{31}=
\frac{1}{\varepsilon^d}\int_{\tau \in
[0,1]}\int_{s\in (-3,-2), (Y,W)\in E} (t-s)^{d+1}(A \nabla_v
f)\left(\tau t+(1-\tau)s,Y,W\right) \\
\times \Bigl[\tau (t-s)\nabla_y\varphi_\var
\Bigl(Y - \tau (t-s)W,\frac{Y-x+(1-\tau)(t-s)W}{\varepsilon}
\Bigr)\\
-\frac{(1-\tau)(t-s)}{\var}\nabla_w\varphi_\var
\Bigl(Y-\tau (t-s)W,\frac{Y-x+(1-\tau)(t-s)W}{\varepsilon}
\Bigr)\Bigr].
\end{multline*}
Using the bounds on the derivatives of $\varphi_\var$ then
yields
\begin{multline}
I_{31}(t,x,v) \lesssim \frac{1}{ \varepsilon^{d+2}}
\int_{\tau \in [0,1]}\int_{s\in
(-3,-2), (Y,W)\in E} |\nabla_v f|\left(\tau
t+(1-\tau)s,Y,W\right) \\
\implies
\label{estim I31}
\fint_{Q_1} I_{31} \lesssim
\frac{1}{ \varepsilon^{d+2}} \int_{Q_3} |\nabla_v f|.
\end{multline}
The result follows from combining \eqref{estim I2},
\eqref{estim I4}, \eqref{estim I1}, \eqref{estim I323}
and~\eqref{estim I31}.
\end{proof}

\begin{remark}
Note that the regularity $W^{\sigma,1}_{x}$ is only used over a
small trajectory that ``noises'' the position variable $x$ in
$Q_1$ with the velocity $w$ in $Q_1^-$, hence allowing to
integrate by parts the diffusion operator using \emph{only} the
variables in $Q_1^-$. Note also that it is possible to get some
$W^{\sigma',1}_{t,x,v}$ regularity in all variable with
$\sigma' \in (0,\sigma)$ small by the same method as in
Lemma~\ref{lem:Kolmogorov}, however such regularity is too weak
to yield any intermediate value estimate alone. Note also that
the gap in time between $Q_1^-$ and~$Q_1$ is used to make sure
the intermediate velocity $\pspfrac{x+\var w-y}{t-s}$ remains
bounded and the various domains of integration remain bounded
along the velocity variable. In fact, the result is false
without such gap, see Remark~\ref{rem:gap}.
\end{remark}

\subsection{Proof of the Intermediate-Value Lemma}
\label{ss:proof-ivl}

In this subsection, we prove that Proposition~\ref{p:HPI lemma}
implies Theorem~\ref{t:IVL}. Take $f$ a sub-solution
to~\eqref{e:main}--\eqref{e:hyp-coef} on $Q_1$ and
satisfying~\eqref{e: hyp IVL} for some given
$\delta_1,\delta_2>0$:\vspace*{-3pt}
\begin{align}
\label{hyp bis IVL}
|\{f\leq 0\}\cap Q_{r_0}^{-}| \geq \delta_1 |Q_{r_0}^{-}|
\quad \text{ and }\ \quad
|\{f\geq 1-\theta\} \cap Q_{r_0}| \geq \delta_2 |Q_{r_0}|.
\end{align}
Define $g:=f-(t+25r_0^2) \|S\|_{L^\infty(Q_1)}$. Then its
positive part $g_+$ is a sub-solution
to~\eqref{e:main}--\eqref{e:hyp-coef} in $Q_{5r_0}$ with zero
source term and with $g_+ \in [0,1]$ since $f\le 1$ in
$Q_{\sfrac{1}{2}}$. We set\vspace*{-3pt}
\[
r_0 = \begin{dcases}
\Bigl(\dfrac{\delta_1}{400
(1+\|S\|_{L^\infty(Q_1)})}\Bigr)^{\sfrac12} \leq
\dfrac{1}{20}&\text{if $S$ non-zero},\\
\dfrac{1}{20}&\text{if $S=0$},
\end{dcases}
\]
and we apply~\eqref{HPI inequality} to $g_+$ at scale $r_0$, for
some $\var>0$ to be chosen later:\vspace*{-3pt}
\begin{equation}\label{Poinc in proof}
\begin{aligned}
\fint_{Q_{r_0}} \bigl(g_+-\langle
g_+\rangle_{Q_{r_0}^{-}}\bigr)_{+}
& \lesssim
\frac{r_0}{\varepsilon^{d+2}} \fint_{Q_{5r_0}} |\nabla_v g_+|
+ \var^\sigma \biggl( \fint_{Q_{5r_0}} g_+^2 \biggr)^{\sfrac12}\\
& \lesssim
\frac{1}{r_0^{4d+1} \varepsilon^{d+2}} \int_{Q_{5r_0}} |\nabla_v g_+|
+ \var^\sigma,
\end{aligned}
\end{equation}
where we have used the bound $g_+ \in [0,1]$ to control the $L^2$
norm. Then~\eqref{hyp bis IVL} implies\vspace*{-3pt}
\begin{equation}
\begin{aligned}
\label{estim averag}
\langle g_+\rangle_{Q_{r_0}^{-}} &= \fint_{(s,y,w) \in
Q_{r_0}^-} \big[f(s,y,w) - (s+25r_0^2) \| S \|_{L^\infty(Q_1)}\big]_+\\
&\leq \frac{\left|{\{f>0\}\cap Q_{r_0}^-}\right|}{|Q^{-}_{r_0}|}
\leq 1-\delta_1
\end{aligned}
\end{equation}
and\vspace*{-3pt}
\begin{equation}\label{lower bd Poinc}
\begin{aligned}
\fint_{Q_{r_0}} \bigl(g_+&-\langle g_+
\rangle_{Q_{r_0}^{-}}\bigr)_{+} \\
&\ge \frac{1}{|Q_{r_0}|} \int_{(t,x,v) \in Q_{r_0}}
\big[ f(t,x,v) - (t+25r_0^2) \|S\|_{L^\infty(Q_1)}
-(1-\delta_1)\big]_{+}\\
&\ge \frac{1}{|Q_{r_0}|} \int_{(t,x,v) \in Q_{r_0}}
\big[ f(t,x,v) - 25r_0^2 \|S\|_{L^\infty(Q_1)}
-(1-\delta_1)\big]_{+}\\
&\geq \frac{1}{|Q_{r_0}|}
\int_{\{f\geq 1-\theta\}\cap Q_{r_0}}
\Bigl(\frac{\delta_1}{2}-\theta\Bigr)_{+}
\ge \delta_2 \Bigl(\frac{\delta_1}{2}-\theta\Bigr).
\end{aligned}
\end{equation}

We then estimate from above the right hand side of the Poincaré
inequality \eqref{Poinc in proof}:\vspace*{-3pt}
\begin{align*}
\int_{Q_{5r_0}} |\nabla_v g_+| &\le \int_{Q_{5r_0}}
|\nabla_v f_+| \\
&\le \int_{\{ f=0 \} \cap Q_{5r_0}} \hspace*{-2mm}\cdots\
+\int_{\{ 0<f<1-\theta \} \cap Q_{5r_0}} \hspace*{-2mm}\cdots\
+ \int_{\{ f\geq 1-\theta \} \cap Q_{5r_0}} \hspace*{-2mm}\cdots\\
&=: I_1 +I_2 + I_3.
\end{align*}
The first term $I_1=0$ since $\nabla_v f_+ = 0$ almost everywhere
on $\{f_+=0\}$
(see~\cite[\S 4.2.2]{MR3409135}). Combining the
Cauchy-Schwarz inequality, Proposition \ref{prop:EE} and the fact
that $f\leq 1$, we get
\begin{align*}
I_2& \leq |\{0<f<1-\theta\}\cap Q_{5r_0}|^{\sfrac12}\biggl(
\fint_{Q_{5r_0}} |\nabla_v f_+|^2 \biggr)^{\sfrac12} \\
& \lesssim
|\{0<f<1-\theta\}\cap Q_{\sfrac12}|^{\sfrac12}
\biggl(\fint_{Q_{\sfrac12}} f_+^2 \biggr)^{\sfrac12}
\lesssim |\{0<f<1-\theta\}\cap Q_{\sfrac12}|^{\sfrac12}
\end{align*}
and (using that $\nabla_v f$ is zero almost everywhere on
$\{ f = \text{cst} \}$, see
again~\cite[\S 4.2.2]{MR3409135})
\begin{align*}
I_3 &= \int_{Q_{5r_0}} \big|\nabla_v
\big[(f-(1-\theta))_+ +(1-\theta) \big]\big|= \int_{Q_{5r_0}}
\big|\nabla_v \big[f-(1-\theta)\big]_+\big| \\
&\lesssim \biggl( \int_{Q_{5r_0}}
\big|\nabla_v \big[f-(1-\theta)\big]_+\big|^2
\biggr)^{\sfrac12} \\
&\lesssim \biggl( \int_{Q_{\sfrac12}}
\big[f-(1-\theta)\big]_+^2
+ \int_{Q_{\sfrac12}}
\big[f-(1-\theta)\big]_+ |S| \biggr)^{\sfrac12}\\
&\lesssim \theta + \theta^{\sfrac12}\|S\|_{L^\infty(Q_1)}
\lesssim \theta^{\sfrac12}\bigl(1+\|S\|_{L^\infty(Q_1)}\bigr),
\end{align*}
where we have used the energy estimate in
Proposition~\ref{prop:EE} on $[f-(1-\theta)]_+$.

The last two estimates on $I_2$ and $I_3$ yield the following
control on the right hand side of~\eqref{Poinc in proof}:
\begin{multline}
\label{upp bd Poinc}
\frac{1}{r_0^{4d+1} \varepsilon^{d+2}}
\int_{Q_{5r_0}} |\nabla_v g_+| +
\varepsilon^\sigma\\
\lesssim
\frac{\theta^{\sfrac12}\left(1+\|S\|_{L^\infty(Q_1)}\right)}{r_0^{4d+1}
\varepsilon^{d+2}}
+ \frac{|\{0<f<1-\theta\}\cap
Q_{\sfrac12}|^{\sfrac12}}{r_0^{4d+1}\varepsilon^{d+2}}
+ \varepsilon^\sigma.
\end{multline}

Combining \eqref{lower bd Poinc} and \eqref{upp bd Poinc} gives,
for some universal constant $C\geq 1$:
\begin{equation}
\label{estim finale}
\frac{\delta_1 \delta_2}{2} \le \delta_2 \theta +
\frac{C\left(1+\|S\|_{L^\infty(Q_1)} \right)
\theta^{\sfrac12}}{r_0^{4d+1} \varepsilon^{d+2}}
+ \frac{C|\{0<f<1-\theta\}\cap
Q_{\sfrac12}|^{\sfrac12}}{r_0^{4d+1} \var^{d+2}} + C \varepsilon^\sigma.
\end{equation}
We choose $\varepsilon$ such that
$C\varepsilon^\sigma \leq \sfrac{\delta_1 \delta_2}{8}$ and
$\theta$ such that
\[
\delta_2 \theta + \frac{C\left(1+\|S\|_{L^\infty(Q_1)}\right)
\theta^{\sfrac12}}{r_0^{4d+1}\var^{d+2}} \leq \frac{\delta_1
\delta_2}{8},
\]
\eg
\begin{equation}
\label{eq:choice-theta}
\varepsilon=\Bigl(\frac{\delta_1\delta_2}{8C}\Bigr)^{\sfrac{1}{\sigma}},
\qquad
\theta= \delta_1^2 \delta_2^2 \biggl[ 8\Bigl( \delta_2 +
\frac{C\left(1+\|S\|_{L^\infty(Q_1)}\right)}{r_0^{4d+1} \left(
\sfrac{\delta_1\delta_2}{8C}
\right)^{\psfrac{d+2}{\sigma}}}\Bigr) \biggr]^{-2},
\end{equation}
which finally implies the result with
\begin{equation}
\label{eq:choice-nu}
\nu:= \frac{1}{|Q_{\sfrac12}|}\biggl(\frac{\delta_1\delta_2}{4C}
\Bigl(\frac{\delta_1\delta_2}{8C}\Bigr)^{\psfrac{d+2}{\sigma}}
r_0^{4d+1}\biggl)^2 \gtrsim
\frac{\left( \delta_1 \delta_2
\right)^{10d+16}}{\left(1+\|S\|_{L^\infty(Q_1)}\right)^{4d+2}}.
\end{equation}

\subsection{Measure-to-pointwise estimate}
\label{ss:m2p}

In this subsection, we combine Proposition~\ref{prop:1st lemma}
and Theorem \ref{t:IVL} to prove a measure-to-pointwise estimate
of ``lowering of the maximum'' à la De~Giorgi.

\begin{lemma}[Measure-to-pointwise upper bound]
\label{l:increase}
Given $\delta \in (0,1)$, define
$r_0 = (\sfrac{\delta}{800})^{\sfrac12}$ if $S$ non-zero and
$r_0\!=\!\sfrac{1}{20}$ if $S\!=\!0$. There is
\hbox{$\mu\!:=\!\mu(\delta)\sim \delta^{2 (1+\delta^{-10d-16})}\!>\!0$} such
that any sub-solution $f$ to~\eqref{e:main}--\eqref{e:hyp-coef}
in $Q_1$ with $S$ such that
$\| S\|_{L^{\infty}(Q_1)}\le \mu$ and so that $f\leq 1$ in
$Q_{\sfrac{1}{2}}$ and
\begin{align}
\label{eq:measure-hyp-pointwise}
\left| \{ f \leq 0 \} \cap Q_{r_0}^- \right| \ge \delta
\left| Q_{r_0}^- \right|
\end{align}
satisfies $f \le 1-\mu$ in $Q_{\sfrac{r_0}{2}}$, with
$Q_{r_0}^{-} := Q_{r_0}(-2r_0^2,0,0) =(-3r_0^2,-2r_0^2] \times
B_{r_0^3} \times B_{r_0}$ (see Figure~\ref{fig:LVI}).
\end{lemma}

\begin{proof}
In view of Proposition~\ref{prop:1st lemma} and the scaling
invariance, there is $\delta'>0$ depending only on $\lambda$
and $\Lambda$ such that for any $r>0$, any sub-solution $f$ on
$Q_{2r}$ so that $\int_{Q_{r}} f_{+}^2 \leq \delta' |Q_{r}|$
satisfies $f\leq \sfrac{1}{2}$ in $Q_{\sfrac{r}{2}}$ (imposing
$C\mu\le \sfrac{1}{4}$ with $C$ the universal constant in the
estimate of Proposition~\ref{prop:1st lemma} used here). Define
then $\nu, \theta>0$ as
in~\eqref{eq:choice-theta}--\eqref{eq:choice-nu} with
$\delta_1=\delta$ and $\delta_2=\delta'$ and a source term
bounded in $L^\infty$ by $1$.

Define $f_k := \theta^{-k}[f-(1-\theta^k)]$ for $k \ge 0$. The
functions $f_{k}$ are sub-solutions to
\eqref{e:main}--\eqref{e:hyp-coef} for all $k \ge 0$ with a
source term of $L^{\infty}$ norm less than $1$ as long as
$k\leq 1+\sfrac{1}{\nu}$ (assuming
$\| S\|_{L^{\infty}(Q_1)}\le \mu$ so that
$\| S\|_{L^{\infty}(Q_1)} \le \theta^{1+\sfrac{1}{\nu}}$). The
sets $\{ 0<f_{k}<1-\theta\}= \{1-\theta^k<f<1-\theta^{k+1}\}$
are disjoints and each $f_k$ satisfies
\eqref{eq:measure-hyp-pointwise}. If
$\int_{Q_{r_0}} (f_k)_{+}^2 \leq \delta' |Q_{r_0}|$ then
$f_k\leq \sfrac{1}{2}$ in $Q_{\sfrac{r_0}{2}}$ so $f\le 1-\mu$
with $\mu = \sfrac{\theta^k}{2}$ which concludes the proof.
Consider $ 1\leq k_0 \leq 1+\nu^{-1}$ such that
$\int_{Q_{r_0}} (f_{k})_{+}^2 > \delta'|Q_{r_0}|$ for any~$k$ such that $0\leq k\leq k_0$. Then for $k$ such that $0\leq k\leq k_0-1$
\begin{align*}
& \left|\left\{ f_{k}\geq 1-\theta \right\} \cap
Q_{r_0}\right| = \left|\{f_{k+1}\geq
0\} \cap Q_{r_0} \right| \geq \int_{Q_{r_0} } (f_{k+1})_{+}^2
> \delta' |Q_{r_0}|, \\
& \left|\left\{f_{k}\leq 0\right\} \cap Q_{r_0}^{-}\right|
\geq \left|\left\{f\leq 0 \right\}\cap Q_{r_0}^{-}\right|
\geq \delta |Q_{r_0}^-|.
\end{align*}
Theorem \ref{t:IVL} for sub-solutions with source term of norm
$L^{\infty}$ less than $1$ then implies, choosing $r_0 =
(\sfrac{\delta}{800})^{\sfrac12}$,
\[
\left| \left\{ 0<f_{k}< 1-\theta \right\}\cap Q_{\sfrac12}\right|
\geq \nu |Q_{\sfrac12}|.
\]
Summing these estimates and using the fact that the sets are
disjoints we have
\[
|Q_{\sfrac12}|\geq \sum_{k=0}^{k_0-1}
\left| \left\{ 0<f_{k}< 1-\theta \right\}\cap
Q_{\sfrac12} \right| \geq k_0 \nu |Q_{\sfrac12}|.
\]
So $k_0 \leq \nu^{-1}$ which ensures that source terms remain
indeed less than one along the iteration, and we deduce
\begin{align*}
f \leq 1-\frac{\theta^{k_0+1}}{2}\leq
1-\frac{\theta^{\psfrac{1+\nu}{\nu}}}{2} \quad \mbox{ in }
Q_{\sfrac{r_0}{2}},
\end{align*}
which yields
$\mu(\delta) := \sfrac{\theta^{1+\sfrac{1}{\nu}}}{2} \sim
\delta^{2 (1+\delta^{-10d-16})}$.
\end{proof}

\section{Harnack inequalities and Hölder continuity}
\label{sec:Harnack}

\subsection{The Harnack inequalities}

To prove the weak Harnack inequality, we first assume $S=0$, and
re-introduce $S$ in the end. Without source term,
$r_0=\sfrac{1}{20}$ can be taken constant in the
measure-to-pointwise estimate. Consider then $h$ non-negative
super-solution to~\eqref{e:main}--\eqref{e:hyp-coef} on $Q_1$ with
$S=0$. The contraposition of Lemma~\ref{l:increase} on the
sub-solution $g := 1 - \sfrac{h}{M}$ then implies for any
$\delta \in (0,1)$ and $M \sim \delta^{-2 (1+\delta^{-10d-16})}$
that
\begin{equation}
\label{eq:scaling-invariant}
\forall \, Q_r(z) \subset Q_1 \text{ with }
Q_{\sfrac{r}{2}}^{+}(z) \subset Q_1, \
\frac{\left| \{ h > M \} \cap Q_{r}(z) \right|}{\left|
Q_{r}(z) \right|} > \delta \implies \inf_{Q_{\sfrac{r}{2}}^+(z)} h \ge 1,
\end{equation}
where
$Q_{\sfrac{r}{2}}^{+}(z) = Q_{\sfrac{r}{2}}(z +(2r^2,2r^2v,0))$,
for $z=(t,x,v)$, is obtained by inverting the operation
$Q_{\sfrac{r}{2}}(z) \to Q_{r}^-(z)$ in Lemma~\ref{l:increase}
(noting that $Q_r^-(z)=Q_r(z-\nobreak(2r^2,2r^2v,0))$). It implies
(inverting the relation $\delta \to M$ and using the layer-cake
representation) that if $\inf_{Q_{\sfrac{r_0}{2}}} h <1$,
\begin{multline}
\label{eq:initial}
\forall \, M \ge 1, \quad \frac{\left| \{ h \ge M \} \cap
Q_{r_0}^- \right|}{\left|Q_{r_0}^- \right|} \lesssim
\delta(M) =\Bigl( \frac{1}{ \ln (1+M)}
\Bigr)^{\spfrac{1}{10d+17}} \\
\implies
\int_{Q_{r_0}^- } \left[ \ln \left( 1+ h \right)
\right]^{\spfrac{1}{10d+18}} \lesssim 1.
\end{multline}
This ``point-to-measure'' estimate controls the decay of the
upper level set in the manner of a \emph{weak Harnack
inequality}, although with a ``logarithmic'' rather than
power-law integrability. We shall now improve the integrability
to a power-law by going back to~\eqref{eq:scaling-invariant} and
performing an inductive argument inspired from the elliptic
theory~\cite{MR3565366}. Note that the logarithmic integrability
in~\eqref{eq:initial} is reminiscent of Moser's approach.

We improve inductively the control of upper level sets in the
following decreasing sequence of cylinders
\begin{equation*}
\cQ^k:= Q_{\Psfrac{r_0}{2}+\alpha_k} \left(-\frac{5}{2}\, r_0^2 +
\frac12 \left( \frac{r_0}{2}+\alpha_k \right)^2,0,0 \right)\subset
Q_{r_0}^{-} \quad \text{ with } \quad
\alpha_k:=\frac{r_0}{2\times 7^{k-1}}.
\end{equation*}
These cylinders satisfy
$\tilde Q^-_{\sfrac{r_0}{2}} \subset \cQ^k \subset \bar \cQ^k
\subset \mathring{\cQ}^{k-1} \subset Q^- _{r_0}$ for all
$k \ge 1$. We now claim that for $\delta_0>0$ small enough (to be
chosen later), for any non-negative super-solution~$h$ with
$\inf_{Q_{\sfrac{r_0}{2}}} h <1$ we have
\begin{align}
\label{eq:claim}
\forall \, k \ge 1, \quad
\frac{\left| \{ h \ge M^k \} \cap \cQ^k \right|}{\left|
\cQ^k \right|} \le \frac{\delta_0}{210^{(4d+2)k}},
\end{align}
where $M \sim \delta^{-2 (1+\delta^{-10d-16})}$ with
$\delta:=\sfrac{\delta_0}{210^{4d+2}}$ as
in~\eqref{eq:scaling-invariant}. Admitting first~\eqref{eq:claim}
we deduce by layer-cake representation that
there is an explicit $\zeta \gtrsim \delta_0^{10d+17}>0$ such
that $\int_{\tilde Q_{\sfrac{r_0}{2}}^{-}} h^\zeta \dd z \lesssim
1$, which implies by linearity
\[
\biggl( \int_{\tilde Q_{\sfrac{r_0}{2}}^{-}} h(z)^\zeta \dd z
\biggr)^{\sfrac{1}{\zeta}} \lesssim \inf_{Q_{\sfrac{r_0}{2}}} h.
\]
This implies the weak Harnack
inequality~\eqref{eq:w-Harnack-stat} on any $f$ non-negative
super-solution to~\eqref{e:main}--\eqref{e:hyp-coef} by applying
the previous estimate to $h := f + (1+t) \|
S\|_{L^\infty(Q_1)}$. To~deduce the Harnack
inequality~\eqref{eq:s-Harnack-stat} we consider $f$ a
non-negative solution to~\eqref{e:main}--\eqref{e:hyp-coef} and
combine the previous control with Proposition~\ref{prop:1st
lemma} to get
\begin{align*}
\sup_{\tilde Q_{\sfrac{r_0}{4}}^{-}} f &\lesssim
\biggl( \int_{\tilde Q_{\sfrac{r_0}{2}}^{-}} f(z)^\zeta \dd z
\biggr)^{\sfrac{1}{\zeta}} + \| S\|_{L^\infty(Q_1)}\\
&\lesssim
\inf_{Q_{\sfrac{r_0}{2}}} f + \| S\|_{L^\infty(Q_1)} \lesssim
\inf_{Q_{\sfrac{r_0}{4}}} f + \| S\|_{L^\infty(Q_1)}.
\end{align*}

Let us now prove the claim~\eqref{eq:claim} to conclude the
proof. The initialization $k=1$ is proved
in~\eqref{eq:initial}. Then define
$A_{k+1} := \{ h > M^{k+1} \} \cap \cQ^{k+1}$ and denote
the following translated centered cylinders
\[
\mathfrak C_{r}[z] :=z\circ Q_{2r}((2r^2,0,0))= z \circ
\left(-2r^2,2r^2\right]\times B_{(2r)^3} \times B_{2r}.
\]
Let us
construct $z_\ell=(t_\ell,x_\ell,v_\ell)\in \cQ^{k+1}$ and
$r_\ell >0$, $\ell \ge 1$, so that:
\begin{enumerate}
\item\label{enum:1} $\forall \, \ell \ge 1$, $r_\ell \in (0,\sfrac{\alpha_{k+1}}{15})$,
\item\label{enum:2} $\forall \, \ell \ge 1$,
$|A_{k+1} \cap \mathfrak C_{15r_\ell}[z_\ell]| \le \delta_0
|\mathfrak C_{15 r_\ell}[z_\ell]|$,
\item\label{enum:3} $\forall \, \ell \ge 1$,
$|A_{k+1} \cap \mathfrak C_{r_\ell}[z_\ell]| > \delta_0
|\mathfrak C_{r_\ell}[z_\ell]|$,
\item\label{enum:4} the cylinders $\mathfrak C_{3r_\ell}[z_\ell]$, $\ell \ge
1$, are disjoint,
\item\label{enum:5} $A_{k+1}$ is covered by the family $\mathfrak
C_{15r_\ell}[z_\ell]$, $\ell \ge 1$.
\end{enumerate}
Note that inverting the operation
$Q_{\sfrac{r}{2}}(z) \to Q_{r}^-(z)$ in Lemma~\ref{l:increase}
yields, when starting from $\mathfrak C_{r_\ell}[z_\ell]$, the
cylinder
$\mathfrak C_{r_\ell}[z_\ell]^+ :=z_\ell\circ
Q_{r_\ell}((10r_\ell^2,0,0))= z_\ell \circ
\left(9r_\ell^2,10r_\ell^2\right]\times B_{r_\ell^3} \times
B_{r_\ell}$. Note also that
$\mathfrak C_{r_\ell}[z_\ell]^+ \subset \mathfrak C_{3
r_\ell}[z_\ell]$ and that property \eqref{enum:1} combined with
$z_\ell \in \cQ^{k+1}$ imply
$\mathfrak C_{15r_\ell}[z_\ell] \subset \cQ^k$. Let us prove
that the family $\mathcal F$ of cylinders $\mathfrak C_{r}[z]$
with $z \in \cQ^{k+1}$, $r \in (0,\sfrac{\alpha_{k+1}}{15})$ and
so that
$|A_{k+1} \cap \mathfrak C_{15r}[z]| \le \delta_0 |\mathfrak
C_{15 r}[z]|$ and
$|A_{k+1} \cap \mathfrak C_{r}[z]| > \delta_0 |\mathfrak
C_{r}[z]|$ cover $A_{k+1}$. We have, using~\eqref{eq:claim} at
the previous step $k$,
\begin{multline}
\label{eq:constraint-varphi}
\forall \, r \in
\left(\frac{\alpha_{k+1}}{15},\alpha_{k+1}\right), \quad
|A_{k+1} \cap \mathfrak C_{r}[z]| \le |A_{k} \cap \mathfrak
C_{r}[z]| \le |A_{k} \cap \cQ^k| \\
\le \frac{\delta_0}{210^{(4d+2)k}} |\cQ^k| \le \delta_0
|\mathfrak C_{r}[z]|.
\end{multline}
If $z \in A_{k+1}$ is not covered by $\mathcal F$ it means that
the continuous positive function
$\varphi(r) = \sfrac{|A_{k+1} \cap \mathfrak C_{r}[z]|}{|\mathfrak
C_{r}[z]|}$ on $(0,+\infty)$ satisfies $\varphi(r) \le \delta_0$
or $\varphi(15r) > \delta_0$ for all
$r \in (0,\sfrac{\alpha_{k+1}}{15})$. The
constraint~\eqref{eq:constraint-varphi} and the continuity impose
$\varphi(r) \le \delta_0$ for all
$r \in\nobreak (0,\sfrac{\alpha_{k+1}}{15})$. Taking $r \to 0$, a
straightforward variation of the Lebesgue differentiation theorem
then implies $z\notin A_{k+1}$ which contradicts the
assumption. Hence $A_{k+1}$ is covered by the family $\mathcal
F$.

It implies in particular that $A_{k+1}$ is covered by the family
$\mathcal F'$ of cylinders $\mathfrak C_{3r}[z]$ with $z \in \cQ^{k+1}$,
$r \in (0,\sfrac{\alpha_{k+1}}{15})$ and such that
$|A_{k+1} \cap \mathfrak C_{15r}[z]| \le \delta_0 |\mathfrak C_{15
r}[z]|$ and
$|A_{k+1} \cap \mathfrak C_{r}[z]| > \delta_0 |\mathfrak
C_{r}[z]|$. The Vitali covering lemma then gives the existence of
a countable sub-family, denoted
$(\mathfrak C_{r_\ell}[z_\ell])_{\ell \ge 1}$, such that the
$(\mathfrak C_{15 r_\ell}[z_\ell])_{\ell \ge 1}$ cover $A_{k+1}$
and the $(\mathfrak C_{3 r_\ell}[z_\ell])_{\ell \ge 1}$ are
disjoint. The Vitali lemma applies thanks to the following
property:
\begin{equation*}
\Big[ \mathfrak C_{r_1}[z_1] \cap \mathfrak C_{r_2}[z_2] \neq
\varnothing \ \mbox{ and } \ r_1\le 2r_2 \Big] \implies \mathfrak C_{r_1}[z_1] \subset
\mathfrak C_{5r_2}[z_2].
\end{equation*}
Take $z_0=(t_0,x_0,v_0)$ in the intersection and
$z=(t,x,v)\in \mathfrak C_{r_1}[z_1]$. Inequalities
$|t-t_2|\leq 18 r_2^2$ and $|v-v_2|\leq 10r_2$ come naturally and
$|x-[x_2+2 r_2^2v_2 +(t-t_2)v_2]| \le 200r_2^3$ follows from
\begin{align*}
\big| x&- \big[ x_2+ 2r_2^2v_2 + (t-t_2)v_2 \big] \big| \\
&\le \big| x- \big[ x_1+ 2r_1^2v_1 + (t-t_1)v_1 \big] \big|\\
&\hspace*{4cm} +
\big| \big[ x_2+ 2r_2^2v_2 + (t-t_2)v_2 \big] - \big[ x_1+
2r_1^2v_1 + (t-t_1)v_1 \big] \big| \\
& \le r_1^3 + \big| \big[ x_2+ 2r_2^2v_2 + (t_0-t_2)v_2 \big] -
\big[ x_1+ 2r_1^2v_1 + (t_0-t_1)v_1 \big] \big|\\
&\hspace*{4cm} + \big|
(t-t_0)(v_2-v_1) \big| \\
& \le 128 r_2^3 + \big| x_0- \big[ x_1+ 2r_1^2v_1 +
(t_0-t_1)v_1 \big] \big| + \big| x_0- \big[ x_2+ 2r_2^2v_2 +
(t_0-t_2)v_2 \big] \big|\\
&\le 200 r_2^3.
\end{align*}
This finishes constructing the covering with the properties
(1)-(2)-(3)-(4)-(5) above. Then Lemma~\ref{l:increase} applied to
each $\mathfrak C_{r_\ell}[z_\ell]$ implies
$\mathfrak C_{r_\ell}[z_\ell]^+ \subset A_k$, and the
$\mathfrak C_{r_\ell}[z_\ell]^+ \subset \mathfrak C_{3
r_\ell}[z_\ell]$ are disjoint. We deduce
\begin{align*}
|A_{k+1}|
& \le \sum_{\ell \ge 1} |A_{k+1} \cap \mathfrak C_{15
r_\ell}[z_\ell]| \le \delta_0 \sum_{\ell \ge 1}
|\mathfrak C_{15 r_\ell}[z_\ell]| \\
& \le 15^{4d+2} \delta_0 \sum_{\ell \ge 1} |\mathfrak
C_{r_\ell}[z_\ell]| \le 30^{4d+2} \delta_0 \sum_{\ell \ge 1}
|\mathfrak C_{r_\ell}[z_\ell]^+| \\
& \le 30^{4d+2} \delta_0 |A_k|
\le \frac{30^{4d+2}\delta_0^2}{210^{(4d+2)k}}
\le \frac{\delta_0}{210^{(4d+2)(k+1)}} \left| \cQ^{k+1}
\right|
\end{align*}
for $\delta_0$ small enough which proves the induction
claim~\eqref{eq:claim} and concludes the proof.

\subsection{The Hölder continuity}
\label{ss:Holder}

De~Giorgi's argument to Hölder continuity uses the
measure-to-pointwise Lemma~\ref{l:increase}. We briefly sketch it
in order to track the constant. Hölder regularity could also be
deduced from the Harnack inequality in
Theorem~\ref{t:harnack}. Given $f$ solution to
\eqref{e:main}--\eqref{e:hyp-coef} on $Q_2$ and
$r_0 =\sfrac{1}{40}$
\begin{equation}
\label{eq:reduc-osc}
\osc_{Q_{r_0}}f \leq \left( 1- \frac{\mu}{2} \right)
\max\bigl(\osc_{Q_{1}}f,e^{2(1+2^{10d+16})} \|S\|_{L^{\infty}(Q_2)}\bigr)
\end{equation}
follows from~Lemma \ref{l:increase} rescaled to $Q_2$ with
$\delta=\sfrac12$ and applied to whichever of $F$ or $-F$
satisfies~\eqref{eq:measure-hyp-pointwise}, where
\[\textstyle
F : =2 \bigl[\max\bigl(\osc_{Q_{1}} f,
e^{2(1+2^{10d+16})}\|S\|_{L^{\infty}(Q_2)} \bigr) \bigr]^{-1} \bigl[f-\frac12
(\sup_{Q_{1}} f+\inf_{Q_{1}} f)\bigr].
\]
By iteration we deduce
\begin{multline}
\label{eq:iter-osc}
\forall z_0 \in Q_{1}, \forall \, r\in(0,r_0),\\
\osc_{Q_{r}(z_0)}f
\leq r^{\alpha} e^{2(1+2^{10d+16})}
\left(1+\|S\|_{L^{\infty}(Q_2)}\right) \max\bigl(e^{2(1+2^{10d+16})}
\|S\|_{L^{\infty}(Q_2)},\osc_{Q_1}f \bigr).
\end{multline}
Indeed, the following sequence of solution of
\eqref{e:main}--\eqref{e:hyp-coef} in $Q_1$
\begin{align*}
f_{n}(\tau,y,w)= 2
\frac{\left( 1 - \sfrac{\mu}{2}
\right)^{1-n} f\left(t_0+ r_0^{2n} \tau, x_0
-r_0^{2n} \tau v_0 + r_0^{3n} y,v_0+ r_0
^n w \right)}{\max \bigl(\osc_{Q_{1}}f, e^{2(1+2^{10d+16})}
\|S\|_{L^{\infty}(Q_2)} \bigr)}
\end{align*}
satisfies
$\osc_{Q_{1}}f_{n} \leq 2
e^{2(1+2^{10d+16})}(1+\|S\|_{L^{\infty}(Q_2)})$ by induction on
$n \!\ge\! 1$ (the case $n\!=\!1$ is true by definition of $f_n$ and it
propagates thanks to~\eqref{eq:reduc-osc}). If one defines
$\alpha\in (0,1)$ such that $1-\sfrac{\mu}{2}=r_0^{\alpha}$
(assuming that $\mu$ is small enough), the previous induction
implies~\eqref{eq:iter-osc} by standard arguments. To deduce the
Hölder estimate between $z,z'\in\nobreak Q_1$, use intermediate
points in $[z,z']$ at distance less than $r_0$ and the
estimate~\eqref{eq:iter-osc}.

\vspace*{\baselineskip}\enlargethispage{-\baselineskip}
\backmatter
\bibliographystyle{jepalpha+eid}
\bibliography{guerand-mouhot}
\end{document}
